% concatenated from HJT_0808V1_.tex
\documentclass[11pt]{amsart}

\usepackage[utf8]{inputenc}
\usepackage{lmodern}
%\usepackage[scrtime]{prelim2e}
\usepackage{amsmath,amssymb,amsthm,latexsym,cite,cancel}
\usepackage[small]{caption}
\usepackage{graphicx,wasysym,overpic,tikz,color}
\usepackage{subfigure}
\usepackage{cite}
\usepackage[colorlinks=true,urlcolor=blue,
citecolor=red,linkcolor=blue,linktocpage,pdfpagelabels,
bookmarksnumbered,bookmarksopen]{hyperref}
\usepackage[english]{babel}
\usepackage{units}
\usepackage[utf8]{inputenc}
\usepackage{enumitem}
\usepackage[left=2.1cm,right=2.1cm,top=2.71cm,bottom=2.71cm]{geometry}
\usepackage[hyperpageref]{backref}
\usepackage{float}
\usepackage[T1]{fontenc}

%\usepackage{refcheck}
%\usepackage{showkeys}
\usepackage[colorinlistoftodos]{todonotes}
%
%\makeatletter
%\providecommand\@dotsep{5}
%\def\listtodoname{List of Todos}
%\def\listoftodos{\@starttoc{tdo}\listtodoname}
%\makeatother
%%%



%\usepackage[latin1]{inputenc}
%\usepackage{amsfonts}
%\usepackage{amscd}
%\usepackage{latexsym}
%\usepackage{mathrsfs}
%\usepackage[notref,notcite]{showkeys}
% \font\rmg=cmr10
%\font\bfg=cmbx10
%\renewcommand{\arraystretch}{1.5}
%\textwidth172mm \textheight22cm \hoffset-24mm \voffset-20mm
%\font\sans=cmss10
%\textwidth=13.5cm
%\baselineskip=16pt
%\makeatletter \@addtoreset{equation}{section} \makeatother
%\renewcommand{\theequation}{\arabic{section}.\arabic{equation}}
%\setlength{\parindent}{1em}

\newtheorem{theorem}{Theorem}[section]
\newtheorem{definition}[theorem]{Definition}
\newtheorem{proposition}[theorem]{Proposition}
\newtheorem{lemma}[theorem]{Lemma}
\newtheorem{remark}[theorem]{Remark}
\newtheorem{example}[theorem]{Example}
\newtheorem{corollary}[theorem]{Corollary}
%\usepackage{CJK,CJKnumb,CJKulem,times,dsfont,ifthen,mathrsfs,latexsym,amsfonts,color}
%\usepackage{amsmath,amsthm,makeidx,fontenc,amssymb,bm,graphicx,psfrag,listings,curves,extarrows}

\newcommand{\abs}[1]{\lvert#1\rvert}
\newcommand{\norm}[1]{\lVert#1\rVert}
\newcommand{\red}[1]{\textcolor{red}{#1}}
\newcommand{\blue}[1]{\textcolor{blue}{#1}}
\definecolor{revisionpurple}{rgb}{0.5,0,0.5}
\newcommand{\purple}[1]{\textcolor{revisionpurple}{#1}}
\newcommand{\D}{\mathcal{D}}
\newcommand{\G}{\mathcal{G}}
\newcommand{\K}{\mathcal{K}}
\newcommand{\h}{\mathcal{H}}
\newcommand{\E}{\mathrm{E}}
\newcommand{\V}{\mathrm{V}}
\newcommand{\R}{\mathbb{R}}
\newcommand{\C}{\mathbb{C}}
\newcommand{\e}{\mathrm{e}}
\newcommand{\vv}{\mathrm{v}}
\newcommand{\w}{\mathrm{w}}
\newcommand{\W}{\mathrm{W}}


\newcommand{\green}[1]{\textcolor{green}{#1}}


\tikzstyle{nodo}=[circle,draw,fill,inner sep=0pt,minimum size=%
1.5mm]
\tikzstyle{bnodo}=[circle,draw,fill,inner sep=0pt,minimum size=%
2mm]

\numberwithin{equation}{section}
%\newtheorem{theorem}{Theorem}[section]
%\newtheorem{proposition}[theorem]{Proposition}
%\newtheorem{lemma}[theorem]{Lemma}
%\newtheorem{remark}[theorem]{Remark}
%\newtheorem{example}[theorem]{Example}
%\newtheorem{corollary}[theorem]{Corollary}
%\newtheorem{definition}[theorem]{Definition}
%\theoremstyle{definition}


\title[The Schr\"odinger equations with nonlinear point defects on noncompact metric graphs]{ Normalized solutions for $L^2$-supercritical Schr\"odinger equations with nonlinear point defects on noncompact metric graphs}
\author[Z. He]{Zhentao He}

\address[Z. He]{\newline\indent
	School of Mathematics
	\newline\indent
	East China University of Science and Technology
	\newline\indent
	Shanghai 200237, PR China }
\email{\href{mailto:hezhentao2001@outlook.com}{hezhentao2001@outlook.com}}

\author[C. Ji]{Chao Ji}

\address[C. Ji]{\newline\indent
	School of Mathematics
	\newline\indent
	East China University of Science and Technology
	\newline\indent
	Shanghai 200237, PR China }
\email{\href{mailto:jichao@ecust.edu.cn}{jichao@ecust.edu.cn}}

\author[Y. Tao]{Yifan Tao}

\address[Y. Tao]{\newline\indent
	School of Mathematics
	\newline\indent
	East China University of Science and Technology
	\newline\indent
	Shanghai 200237, PR China }
\email{\href{mailto:taoyf2002@126.com}{taoyf2002@126.com}}

\subjclass[2020]{35R02, 81Q35, 58E05, 35Q40}
%
\date{\today}
\keywords{Nonlinear Schr\"odinger equations, $\delta$-type interactions, $L^2$-supercritical, Noncompact metric graphs, Multiplicity}


\begin{document}
\begin{abstract}
In this paper, we study the existence and multiplicity of normalized solutions for the following $L^2$-supercritical Schr\"odinger equation with nonlinear point defects on a noncompact metric graph $\G=(\V,\E)$  
\begin{equation*}
        \begin{cases} 
u'' = \lambda u & \text{on every } \e \in \E, \\ 
\int_\G \abs{u}^2\, dx = \mu & \\
\displaystyle \sum_{\e \succ \vv}  u'_\e(\vv) = -|u(\vv)|^{p-2}u(\vv) & \text{at every } \vv \in \V_0, \\ 
\displaystyle \sum_{\e \succ \vv}  u'_\e(\vv) = 0 & \text{at every } \vv \in \V \setminus \V_0,
\end{cases}
\end{equation*}
where $p>4$,  $\G$ has finitely many edges and no self-loops, $\V_0 \subset \V$ is nonempty, $\mu>0$ is a given constant, $\lambda$ is an unknown Lagrange multiplier, $\e \succ \vv$ means that the edge $\e$ is incident at $\vv$, and the notation $u'_\e(\vv)$ stands for $u'_\e(0)$ or $-u'_\e(\ell_\e)$, according to whether the vertex $\vv$ is identified with $0$ or $\ell_\e$.  We first prove the existence of a positive normalized solution for every prescribed mass and every nonempty set of defect vertices. We then establish a multiplicity result for normalized solutions when the prescribed mass is sufficiently small and sufficiently many half-lines are attached to the defect vertices.
\end{abstract}
\maketitle


\section{Introduction}
In this paper,  we consider a connected metric graph $\G = (\V, \E)$, where $\E$ is the set of edges and $\V$ is the set of vertices. Each bounded edge $\e$ is identified with a closed and bounded interval $I_\e = [0, \ell_\e]$ with $\ell_\e > 0$, while each unbounded edge $\e$ is identified with a copy of $I_\e = \mathbb{R}^+ = [0, +\infty)$ and is called a half-line.
%If $\G$ is a metric graph with a finite number of edges, its compact core $\K$ is defined as the metric subgraph of $\G$ consisting of all the bounded edges of $\G$. Therefore, if $\G$ is compact, then its compact core $\K = \G$. 

Throughout the paper, we assume that $\mathcal{G}$ has finitely many edges and no self-loops (see, e.g., Figure \hyperref[figure1]{1}). This entails no loss for the present results, since any self-loop can be subdivided at an interior point by inserting a degree-two Kirchhoff vertex outside $\V_0$.
%, and the compact core $\K$ is non-empty, unless otherwise specified. 
A connected metric graph is naturally a locally compact metric space, and for any two points $x,y \in \G$, the distance $\operatorname{dist}(x,y)$ is the length of a shortest path along the edges. Moreover, a connected metric graph with a finite number of vertices is compact if and only if it contains no half-lines.

 Quantum graphs (metric graphs equipped with differential operators) arise naturally as simplified models in mathematics, physics, chemistry, and engineering when one considers the propagation of waves of various types through a quasi-one-dimensional (e.g., mesoscale or nanoscale) system modeled by a thin neighborhood of a graph. For further details on quantum graphs, see, for example, \cite{Be}.
\begin{figure}[H]
	\begin{tikzpicture}[xscale=0.4, yscale=0.4]
		\draw (-14,3)--(-11,6);
		\draw (0,8)--(-8,6);
		\draw (-11,6)--(-8,6);
		\draw (-22,7)--(-14,3);
		\draw (0,3)--(-8,3);
		\draw (-11,9)--(-11,6);
		\draw (-14,3)--(-14,6);
		\draw (-11,3)--(-11,6);
		\draw (-8,3)--(-8,6);

		\node at (-14,3) [nodo] {};
		\node at (-14,6) [nodo] {};
		\node at (-11,3) [nodo] {};
		\node at (-11,6) [nodo] {};
		\node at (-8,3) [nodo] {};
		\node at (-8,6) [nodo] {};
		\node at (-11,9) [nodo] {};

		\node at (-22,7.5) {$\infty$};
		\node at (0,8.5) {$\infty$};
		\node at (0,3.5) {$\infty$};

		\draw (-11,10.5) circle (1.5);
		\node at (-11,12) [nodo] {};
	\end{tikzpicture}
	\caption{A noncompact metric graph $\G$ with finitely many edges, no self-loops and a nonempty compact core.}
	\label{figure1}
\end{figure}
Extensive research has been devoted to the existence and multiplicity of normalized solutions to the following nonlinear Schr\"{o}dinger equation on a metric graph $\G=(\V,\E)$
\begin{equation}
\label{eqs}
    -u'' - \lambda u= \abs{u}^{p-2}u\quad \text{on every } \e \in \E.
\end{equation}
In the $L^2$-subcritical ($2<p<6$) and $L^2$-critical ($p=6$) cases, the energy functional associated with equation \eqref{eqs} is bounded from below and coercive under the mass constraint (requiring the mass to be below a certain threshold when $p=6$). This structure allows for the use of minimization methods to obtain normalized ground states, i.e., solutions with minimal energy under the mass constraint. The existence of such normalized ground states has been established both on compact metric graphs \cite{CDS, Do} and on noncompact metric graphs \cite{ABD, ACFN, AST, AST2, AST3, Do1, NP, PS}. In addition, the existence of local minimizers was studied by Pierotti, Soave and Verzini in \cite{PSV}. For a noncompact metric graph $\G$ with finitely many edges and a nonempty compact core $\K$, the following modification of equation \eqref{eqs}, in which the nonlinearity acts only on $\K$, was introduced in \cite{Gn,No}
\begin{equation}
\label{eqsloc}
    -u'' - \lambda u= \chi_\K\abs{u}^{p-2}u \quad \text{on every } \e \in \E,
\end{equation}
where $\chi_\K$ is the characteristic function of the compact core $\K$. For the existence and nonexistence of normalized ground states for equation \eqref{eqsloc} in the $L^2$-subcritical case, we refer to \cite{ST, T}; for the $L^2$-critical case, see \cite{DT, DT2}. Moreover, in the $L^2$-subcritical case, Serra and Tentarelli \cite{ST1}, by means of genus theory, proved the existence of multiple bound states with negative energy. 

In the $L^2$-supercritical case ($p>6$), the energy functional is no longer bounded from below under the mass constraint. Moreover, the scaling technique---commonly used in the analysis in $\mathbb{R}^N$---is no longer applicable, and no corresponding Pohozaev identity is available, which creates further complications in proving the boundedness of Palais--Smale sequences. In addition, the topology of metric graphs further complicates the problem. To address these challenges, several works have focused on the $L^2$-supercritical case. In \cite{Bort0}, Borthwick, Chang, Jeanjean and Soave proved the existence of bounded Palais--Smale sequences carrying Morse index type information for functionals with minimax geometry under the mass constraint. In \cite{Cha}, Chang, Jeanjean, and Soave studied the nonlinear Schr\"{o}dinger equation \eqref{eqs} on compact metric graphs. By combining blow-up analysis, monotonicity methods, and an analysis of Morse index type information for Palais--Smale sequences and bound states, the authors proved the existence of normalized nonconstant bound states for equation \eqref{eqs} when the prescribed mass is sufficiently small. Subsequently, in \cite{Bort}, for arbitrary prescribed mass, Borthwick, Chang, Jeanjean and Soave investigated the existence of normalized solutions for \eqref{eqsloc} on noncompact metric graphs with finitely many edges and a nonempty compact core $\K$. The multiplicity of normalized solutions to equation \eqref{eqsloc} was proved in \cite{Ca} by employing a more general blow-up analysis. In \cite{Do0}, Dovetta, Jeanjean and Serra studied the existence of normalized solutions for the nonlinear Schr\"{o}dinger equation \eqref{eqs} on periodic metric graphs, as well as on noncompact metric graphs with finitely many edges under suitable topological assumptions.

Very recently, Boni, Dovetta, and Serra \cite{BDS} were the first to investigate the existence of normalized ground states for the following Schr\"odinger equation on a metric graph $\G=(\V,\E)$ with nonlinear point defects
\begin{equation} \label{eqpdbds}
    \begin{cases} 
u'' = \lambda u & \text{on every } \e \in \E, \\ 
\int_\G \abs{u}^2\, dx = \mu & \\
\displaystyle \sum_{\e \succ \vv}  u'_\e(\vv) = -|u(\vv)|^{p-2}u(\vv) & \text{at every } \vv \in \V_0, \\ 
\displaystyle \sum_{\e \succ \vv}  u'_\e(\vv) = 0 & \text{at every } \vv \in \V \setminus \V_0,
\end{cases}
\end{equation}
where $\V_0\subset \V$ is nonempty, $\mu>0$ is prescribed, $p \in (2,4)$ (the $L^2$-subcritical case for equation \eqref{eqpdbds}), and $\lambda \in \mathbb{R}$ is an unknown Lagrange multiplier. For metric graphs with finitely many vertices, the authors proved the existence of ground states for every $\mu > 0$ and $p \in (2,4)$.  For metric graphs with infinitely many vertices, focusing on $\mathbb{Z}$-periodic graphs and the two-dimensional square grid, the authors proved a series of results unraveling nontrivial threshold phenomena for both $\mu > 0$ and $p \in (2,4)$. Their study highlights that  the ground state problem depends critically on the degree of periodicity of the metric graph, on the total number of point defects, and on their dislocation on the metric graphs. For further results concerning the Schr\"{o}dinger equations with nonlinear point defects on metric graphs or in dimension one, we refer the reader to \cite{BD,BD2,BBDT,BBDT2} and the references therein.


The condition
$$
\sum_{\e \succ \vv} u'_\e(\vv) = 0,\quad \forall \vv \in \V \setminus \V_0
$$
is referred to as the Kirchhoff or natural boundary condition, while
$$
\sum_{\e \succ \vv} u'_\e(\vv) = -|u(\vv)|^{p-2}u(\vv), \quad \forall \vv \in \V_0,
$$
can be interpreted as representing the effect of a deep attractive potential well or a strong attractive defect at each vertex $\vv \in \V_0$. In the literature, such a condition is commonly termed a nonlinear $\delta$-interaction and is regarded as modeling strongly localized, point-like defects or inhomogeneities in the propagation medium (see, e.g., \cite{ABR} for a comprehensive overview of non-Kirchhoff vertex conditions on graphs). Models with concentrated nonlinearities have been proposed, for instance, in semiconductor theory \cite{JPS, Ni}, to describe quantum dynamics in resonant tunneling diodes, as well as the effect of the confinement of charges in small regions.

Formally, solutions to problem \eqref{eqpdbds} correspond to standing waves of the following nonlinear Schr\"odinger equation
\[
i\partial_t \psi + \partial_{xx} \psi + \sum_{\vv \in \V_0} \delta_\vv |\psi|^{p-2} \psi = 0 \quad \text{on } \mathcal{G}
\]
via the ansatz $\psi(t,x) = e^{i\lambda t} u(x)$. The constraint on the $L^2$-norm of $u$ in \eqref{eqpdbds},  which reflects the normalization requirement imposed on the solutions, is a standard condition on the mass of $u$ (sometimes interpreted as the number of particles in the condensate)---a quantity that remains preserved along the evolution in time.

We note that \cite{BDS}  studied the existence of normalized solutions to equation \eqref{eqpdbds} only in the  $L^2$-subcritical  case ($2<p<4$). It is therefore natural to investigate the existence of normalized solutions to equation \eqref{eqpdbds} in the  $L^2$-supercritical case ($p>4$). To the best of our knowledge, the mass-supercritical problem has not yet been studied on noncompact metric graphs with finitely many edges and an arbitrary nonempty set of nonlinear point defects. In this paper, we take a first step toward filling this gap. As in the cases of equations \eqref{eqs} and \eqref{eqsloc}, by \cite[Lemma 2.2]{BDS}, in the $L^2$-subcritical ($2<p<4$) and $L^2$-critical ($p=4$) cases, the energy functional associated with \eqref{eqpd} is bounded from below and coercive under the mass constraint (requiring the mass to be below a certain threshold when $p=4$) and one can use minimization methods to obtain normalized ground states. In contrast, in the $L^2$-supercritical case ($p>4$), the energy functional is no longer bounded from  below on the mass constraint (one can see this with the functions $\{\varphi^t\}$ defined in Lemma \ref{lem3.1}). Moreover, due to the lack of a suitable mass-preserving scaling and a corresponding Pohoz\v{a}ev identity for \eqref{eqpd} on general metric graphs, establishing the boundedness of constrained Palais--Smale sequences for the energy functional is particularly challenging.

Motivated by \cite{BDS,Bort,Ca}, in this paper, we study the existence and multiplicity of normalized solutions to the following $L^2$-supercritical Schr\"odinger equation with nonlinear point defects on a noncompact metric graph $\G=(\V,\E)$ 
\begin{equation}\label{eqpd}
        \begin{cases} 
u'' = \lambda u & \text{on every } \e \in \E, \\ 
\int_\G \abs{u}^2\, dx = \mu & \\
\displaystyle \sum_{\e \succ \vv}  u'_\e(\vv) = -|u(\vv)|^{p-2}u(\vv) & \text{at every } \vv \in \V_0, \\ 
\displaystyle \sum_{\e \succ \vv}  u'_\e(\vv) = 0 & \text{at every } \vv \in \V \setminus \V_0,
\end{cases}
\end{equation}
where $p>4$, $\G$ has finitely many edges, $\V_0 \subset \V$ is nonempty, $\mu>0$ is the prescribed mass, $\lambda$ is an unknown Lagrange multiplier, $\e \succ \vv$ means that the edge $\e$ is incident to $\vv$, and the notation $u'_\e(\vv)$ stands for $u'_\e(0)$ or $-u'_\e(\ell_\e)$, according to whether the vertex $\vv$ is identified with $0$ or $\ell_\e$. 

Accordingly, a function $u: \G \to \mathbb{R}$ is actually a family of functions $(u_\e)$, where $u_\e: I_\e \to \mathbb{R}$ is the restriction of $u$ to the edge $\e$. The usual $L^p$ space on the metric graph  $\G$ is defined as follows\blue{:}
$$
L^p(\mathcal{G}):=\bigoplus_{\e \in \mathrm{E}} L^p(I_\e),
$$
endowed with the norm
$$
\norm{u}_{L^p(\G)}^p := \sum_{\e \in \mathrm{E}}\norm{u_\e}_{L^p(I_\e)}^p, \quad \text { if } p \in[1, \infty) \quad \text { and } \quad \norm{u}_{L^{\infty}(\G)}:=\max_{\e \in \mathrm{E}}\norm{u_\e}_{L^{\infty}(I_\e)}.
$$
The inner product in $L^2(\mathcal{G})$ is
\begin{equation*}
(u,v)_2= \int_\G uv\,dx.
\end{equation*}
If $u:\G \to\mathbb{R}$ is continuous on $\G$, then we write $u \in C(\G)$. The Sobolev space $H^1(\G)$ consists of all continuous functions $u : \G \to \mathbb{R}$ such that $u_\e \in  H^1([0,\ell_\e])$ for every edge $\e\in \E$ (where $\ell_\e$ may equal $+\infty$); the norm in $H^1(\G)$ is defined as
\begin{equation*}
    \norm{u}^2_{H^1(\G)} = \underset{\e \in \E}\sum \left(\norm{u'}^2_{L^2(I_\e)}+\norm{u}^2_{L^2(I_\e)}\right),
\end{equation*}
and the inner product in this space is
\begin{equation*}
(u,v)= \underset{\e \in \E}\sum \int_{I_\e} \left(u'_{\e}v'_{\e}+  u_{\e}v_{\e}\right)\,dx.
\end{equation*}
For convenience, we will abbreviate $\norm{u}_{L^p(\G)}$ as $\norm{u}_p$ and $\norm{u}_{H^1(\G)}$ as $\norm{u}$. 

We denote by $E( \cdot,\G): H^1(\G) \to \mathbb{R}$ the energy functional associated with equation \eqref{eqpd}:
\begin{equation}\label{func}
E(u,\G) := \frac{1}{2}\norm{u'}^{2}_2-\frac{1}{p}\sum_{\vv \in \V_0}\abs{u(\vv)}^p.
\end{equation}
We shall study the existence and multiplicity of critical points of the functional $E( \cdot,\G)$ restricted to the $L^2$-sphere
$$
H^1_\mu(\G):=\{u \in H^1(\G) \mid \norm{u}^2_2=\mu\},
$$
where $\mu > 0$.

Our first main result is stated as follows:
\begin{theorem}\label{th1}
     Let $\G$ be any noncompact metric graph with finitely many edges and $p>4$. Then, for every nonempty $\V_0 \subset \V$ and every $\mu>0$, problem \eqref{eqpd} has a solution $(u,\lambda)\in H^1(\G) \times (0,+\infty)$ with $u>0$ on $\G$.
\end{theorem}

In \cite{Bort},  by using a minimax principle for some constrained functionals which combines the monotonicity trick and  “approximate Morse index” information on the Palais-Smale sequences (see \cite{Bort0}), Borthwick, Chang, Jeanjean and Soave first proved the existence of normalized solutions $(u_\rho,\lambda_\rho)$ to the following parameterized nonlinear Schr\"{o}dinger equations
\begin{equation}
\label{eqslocrho}
    \begin{cases}
    -u'' - \lambda u= \rho\chi_\K\abs{u}^{p-2}u & \text{on every } \e \in \E, \\ 
    \int_\G \abs{u}^2\, dx = \mu &
    \end{cases}
\end{equation}
 with $u_\rho>0$ for almost every $\rho \in [\frac{1}{2},1]$. By choosing a sequence $\rho_n\to1^-$ and applying a blow-up analysis to $\{u_{\rho_n}\}$ with bounded Morse index, the authors proved the existence of normalized solutions to \eqref{eqsloc}.
 
However, since the blow-up analysis in \cite{Bort}  relies on the structure of \eqref{eqsloc} (or \eqref{eqs} in \cite{Cha}), it is not directly applicable to equation \eqref{eqpd} due to the presence of point defects. To overcome this difficulty, we employ the theory of ordinary differential equations to establish the boundedness of Lagrange multipliers $\{\lambda_{\rho_n}\}$, where $\{(u_{\rho},\lambda_{\rho})\}$ are normalized solutions to the following parameterized equations
\begin{equation}\label{eqpdrho}
       \begin{cases} 
u'' = \lambda u & \text{on every }\e \in \E, \\ 
\int_\G\abs{u}^2\, dx = \mu & \\
\displaystyle\sum_{\e \succ \vv} u'_\e(\vv) = -\rho|u(\vv)|^{p-2}u(\vv) & \text{at every } \vv \in \V_0,\\
\displaystyle\sum_{\e \succ \vv} u'_\e(\vv) = 0 & \text{at every } \vv \in \V \setminus \V_0,
\end{cases} 
\end{equation}
for almost every $\rho \in [\frac{1}{2},1]$, and choose a sequence $\rho_n\to1^-$ as $n \to \infty$. Indeed, if $\lambda_{\rho_n} \to +\infty$ as $n \to \infty$, then, the boundedness of $\left\{\frac{\sum_{\vv \in \V_0}\abs{u_{\rho_{n}}(\vv)}^p}{\lambda_{\rho_{n}}}\right\}$ leads to a contradiction with $\int_\G\abs{u_{\rho_{n}}}^2\, dx \to 0$ as $n \to \infty$. Using this boundedness result, we can then prove the existence of normalized solutions to equation \eqref{eqpd} via the minimax principle developed in \cite{Bort0}, together with an analysis of Morse index type information for Palais--Smale sequences as in \cite{Cha,Bort}. 

We also note that the blow-up analysis developed in \cite{Bort,Cha} relies on the positivity of $u_\rho$ (see also \cite{Ca} for a more general blow-up analysis that avoids this assumption), where $(u_\rho,\lambda_\rho)$ is a normalized solution to \eqref{eqslocrho}. By contrast, our approach does not require $u_\rho>0$ for solutions $(u_\rho,\lambda_\rho)$ of \eqref{eqpdrho}. This  motivates further investigation into the multiplicity of normalized solutions for equation \eqref{eqpd}.

The contributions of the present paper are twofold. First, we treat the mass-supercritical regime $p>4$ for an arbitrary nonempty set of nonlinear point defects on every finite-edge noncompact metric graph and for every prescribed mass. The main compactness difficulty is the uniform control of the Lagrange multipliers. Since the blow-up arguments developed for nonlinearities distributed over the compact core are not directly available for point interactions, we derive an ODE-based estimate that does not require positivity of the approximating solutions. Second, the nonlinear part of the energy depends only on the finite-dimensional trace vector at the defect vertices. Consequently, the infinite sequence of linking structures available for compact-core nonlinearities cannot be reproduced by the same argument. We overcome this obstruction on a suitable class of graphs by attaching $M^i$ half-lines at the $i$-th defect vertex and deriving vertex-dependent trace estimates. These estimates separate $q$ minimax windows and yield $q$ distinct pairs of normalized solutions for sufficiently small masses. The finite-edge assumption is used in the ODE-based compactness argument, in particular when the edgewise estimates are summed over the graph, and in the finite-dimensional trace construction underlying the multiplicity result. Extending the analysis to locally finite graphs with infinitely many vertices remains open.


Let $\G=(\V,\E)$ be a connected metric graph with finitely many edges, either compact or noncompact, and let
$
\V_0:=\{\vv_1,\ldots,\vv_q\}\subset \V
$
be a nonempty set of vertices.
For any integer $M\geq 2$, let $\G_M=(\V,\E_M)$ be the graph obtained from $\G$ by attaching $M^{i}$ new half-lines at $\vv_i$, $i=1,\ldots,q$. We now present our multiplicity result for solutions of the following problem
\begin{equation}\label{eqpdm}
\begin{cases} 
u'' = \lambda u & \text{on every }\e \in \E_M, \\ 
\int_{\G_{M}}\abs{u }^2\, dx = \mu & \\
\displaystyle\sum_{\e \succ \vv} u_{\e}' (\vv)= -|u(\vv)|^{p-2}u(\vv)  & \text{at every } \vv \in \V_0,\\
\displaystyle\sum_{\e \succ \vv}u_{\e}' (\vv)=0  & \text{at every } \vv \in \V \setminus \V_0,
\end{cases}
\end{equation}
\begin{theorem}\label{th2}
Let $\G_M$ be defined as above and $p>4$. Then there exist $M_0\geq2$ and $\mu_0>0$ such that, for every integer $M\geq M_0$ and every $0<\mu<\mu_0$, problem \eqref{eqpdm} on $\G_M$ has at least $q$ distinct pairs of solutions $\{(\pm u_N,\lambda_N)\}_{N=2}^{q+1} \subset H^1(\G_{M}) \times (0,+\infty)$.
\end{theorem} 
The finite-dimensional nature of the multiplicity argument can be stated precisely. The nonlinear defect term depends only on the trace vector
\[
T(u):=(u(\vv_1),\ldots,u(\vv_q))\in\mathbb R^q.
\]
Let \((u,\lambda)\) be a constrained critical point of \(E_\rho\) with \(\lambda\geq0\), and consider on the tangent space the quadratic form
\[
Q_{u,\rho}(\varphi)
:=\|\varphi'\|_2^2+\lambda\|\varphi\|_2^2
-(p-1)\rho\sum_{i=1}^q|u(\vv_i)|^{p-2}|\varphi(\vv_i)|^2.
\]
If \(T(\varphi)=0\), then
\[
Q_{u,\rho}(\varphi)=\|\varphi'\|_2^2+\lambda\|\varphi\|_2^2\geq0.
\]
Consequently, the restriction of \(T\) is injective on every subspace on which \(Q_{u,\rho}\) is negative definite, and the constrained Morse index is at most \(q\). Thus the trace-based linking construction used here has at most \(q\) independent negative directions. This is a limitation of the present minimax argument, not a nonexistence statement for additional critical points; compare \cite{Ca}.

Another difficulty is to show that these finitely many levels are
distinct. In order to separate the corresponding critical levels, we attach suitable
numbers of half-lines to the defect vertices.  This modifies the estimates for concentrated nonlinearities, see Lemma \ref{lem-trace-K}, and gives
\begin{equation*}
 \frac{(p-4)2^{\frac{3p-4}{4-p}}}{p}M^{\frac{2p(N-1)}{p-4}}\mu^{\frac{p}{4-p}}\le c_{\rho,\mu}^N
\leq
\frac{3L_0^2}{2\ell^2}K^{\frac{2p}{p-4}}_{M,N-1}
\mu^\frac{p}{4-p}, \quad \rho\in [1/2,1],
\end{equation*} see \eqref{ep-levels}.
By choosing the integer \(M\) sufficiently large, these intervals are
pairwise disjoint, and the corresponding critical points have different energies. This
gives \(q\) distinct pairs of normalized solutions.

\iffalse
In addition to the difficulty mentioned above regarding the reliance of the blow-up analysis in \cite{Cha} on the positivity of $u_\rho$ \blue{and how to distinguish the finitely many solutions}, another main challenge in proving the multiplicity result in \cite{Ca} is to exclude the existence of nontrivial solutions to
\begin{equation}
\label{eqslocrhonol}
    -u'' = \rho\chi_\K\abs{u}^{p-2}u, \quad \text{on every } \e \in \E, 
\end{equation}
where $\rho \in [\frac{1}{2},1]$, and $u \in H^1(\G)$ satisfies $\int_{\G} \abs{u}^2\,dx \leq \mu$ and $\frac{1}{2}\int_{\G} \abs{u'}^2\,dx-\frac{\rho}{p}\int_{\K} \abs{u}^p\,dx$ is sufficiently large. By developing a Poho\v{z}aev type identity on metric graphs and employing the theory of ordinary differential equations, Carrillo, Galant, Jeanjean and Troestler \cite{Ca} proved that, for any sequence $\{u_n\} \subset H^1(\G)$ such that $u_n$ is a solution to \eqref{eqslocrhonol} with $\rho=\rho_n\in [\frac{1}{2},1]$, if 
$$\frac{1}{2}\int_{\G} \abs{u_n'}^2\,dx-\frac{\rho_n}{p}\int_{\K} \abs{u_n}^p\,dx \to +\infty,\,\,\,\text{as}\,\,\, n \to \infty,$$
then $\norm{u_n}_2 \to \infty$ as $n \to \infty$ (see \cite[Theorem 3.6]{Bort0}), and therefore there exists no solution $u$ of equation \eqref{eqslocrhonol} such that $\frac{1}{2}\int_{\G} \abs{u'}^2\,dx-\frac{\rho}{p}\int_{\K} \abs{u}^p\,dx$ is large enough and $\int_{\G} \abs{u}^2\,dx \leq \mu$.  Then, with an argument involving approximate Morse index, the authors proved the convergence of Palais-Smale sequences (deduced by \cite[Theorem 1.12]{Bort0}) of energy functional associted with \eqref{eqslocrho} at levels sufficiently high. In contrast to the approach in \cite{Ca},  we establish a connection between equation \eqref{eqpdrho} with $\lambda=0$ and nonlinear Schr\"odinger equations on combinatorial graphs(see Appendix), and we prove that for any $\rho \in [\frac{1}{2},1]$ and any solution $u\in H^1(\G)$ to 
\blue{
\begin{equation}\label{eqpdrhonol}
       \begin{cases} 
u'' = 0 & \text{on every }\e \in \E, \\ 
\displaystyle\sum_{\e \succ \vv} u'_\e(\vv) = -\rho|u(\vv)|^{p-2}u(\vv) & \text{at every } \vv \in \V_0,\\
\displaystyle\sum_{\e \succ \vv} u'_\e(\vv) =0 & \text{at every } \vv \in \V \setminus \V_0,\\
\end{cases} 
\end{equation}
}
there exists a constant  \blue{$S_{p,\K,\vv_\mathcal{H}} > 0$}, depending only on \blue{$p$, $\K$ and $\vv_\mathcal{H}$}, such that
$$\frac{1}{2}\norm{u'}^{2}_2-\blue{\frac{\rho}{p}\sum_{\vv \in \V_0}\abs{u(\vv)}^p} \leq  \blue{S_{p,\K,\vv_\mathcal{H}}}.$$
Consequently, there is no solution of \eqref{eqpdrhonol} for which  
$$\frac{1}{2}\norm{u'}^{2}_2-\blue{\frac{\rho}{p}\sum_{\vv \in \V_0}\abs{u(\vv)}^p} \blue{\geq S_{p,\K,\vv_\mathcal{H}}+1},\,\,\,\text{for any}\,\,\, \rho \in [\frac{1}{2},1].$$
Note that the restriction on $L^2$-norms is crucial for the non-existence result of equation \eqref{eqslocrhonol}. However, for the non-existence result of equation \eqref{eqpdrhonol},  the restriction on $L^2$-norms is not required.
\fi



In Theorems \ref{th1} and \ref{th2}, we obtain several results for problem \eqref{eqpd} on noncompact metric graphs, for which the embedding $H^1(\G) \hookrightarrow L^2(\G)$ is not compact.  In contrast, on compact metric graphs, this embedding is compact, and therefore the perturbation method developed in \cite[Section 5]{AHJ} can be applied to problem \eqref{eqpd} on compact metric graphs with only minor modifications.  For completeness, we also record the following result for problem \eqref{eqpd} on compact graphs with $p>4$.

\begin{theorem}
    Let $\G=(\V,\E)$ be a compact metric graph, let $\V_0\subset\V$, and assume $p>4$. There exists $\mu^*>0$ depending only on $p$ and $\G$ such that, for any $0<\mu<\mu^*$, problem \eqref{eqpd}
has a solution $(u,\lambda) \in H^{1}(\G)\times (-\infty,1)$. Moreover, for any $m \in \mathbb{N}^+$, there exists $\mu_{m}> 0$,  depending on $p$ and $\G$, such that, for any $0 < \mu < \mu_{m}$,  problem \eqref{eqpd} has at least $m$  solutions $(u_1,\lambda_1), (u_2,\lambda_2),\cdots,(u_{m},\lambda_m) \in H^{1}(\G)\times(-\infty,1)$.
\end{theorem}



The rest of the paper is organized as follows. In Section \ref{sect2}, we present some useful notation and preliminary lemmas. Section \ref{sec3} is devoted to proving  the compactness of bounded Palais-Smale sequences. In Section \ref{sect4}, we provide the proof of Theorem \ref{th1}, and Section \ref{sect5} contains the proof of Theorem \ref{th2}.
\section{Preliminaries}\label{sect2}
\iffalse
\begin{equation}\label{eqpd}
    \begin{cases} 
u'' = \lambda u & \text{on every }\e \in \E, \\ 
\int_\G\abs{u}^2\, dx = \mu & \\
\displaystyle\sum_{\e \succ \vv} u'_\e(\vv) = -|u(\vv)|^{p-2}u(\vv) & \text{at every } \vv \in \V,
\end{cases}
\end{equation}

\begin{align}\label{Prho}
\begin{cases} 
(u_{\rho}^N)'' = \lambda_{\rho}^N u_{\rho}^N & \text{on every }\e \in \E, \\ 
\int_\G\abs{u_{\rho}^N }^2\, dx = \mu & \\
\displaystyle\sum_{\e \succ \vv} (u_{\rho}^N)' (\vv)= -\rho|u_{\rho}^N(\vv)|^{p-2}u_{\rho}^N(\vv)  & \text{at every } \vv \in \V,
\end{cases}
\end{align}
\fi
 

In this section, we recall a Gagliardo-Nirenberg inequality relevant to equation \eqref{eqpd}, introduce some basic notation and results from \cite{Bort0,Ca}, and present a key lemma that allows one to recover a solution to \eqref{eqpd} from a sequence of solutions to the parameterized equations \eqref{eqpdrho}.
\subsection{The Gagliardo-Nirenberg inequality}
Let  $\G=(\V,\E)$ be a noncompact metric graph with finitely many edges and $p \geq 2$. By \cite[Section 2]{AST2}, the following Gagliardo-Nirenberg inequality holds
\begin{equation}\label{eqgn}
        \sum_{\vv \in \V} \abs{u(\vv)}^p \leq C_{p,\G} \norm{u'}_2^{\frac{p}{2}}\norm{u}_2^{\frac{p}{2}}, \quad \forall u \in H^1(\G),
\end{equation}
where $C_{p,\G}:=2^\frac{p}{2}\#\V$.

\subsection{Basic notation and results}\label{subsectnotation}
We first recall some basic notation and results from \cite{Bort0,Bort}.
Let $(E, \langle \cdot, \cdot \rangle)$ and $(H, (\cdot, \cdot))$ be  infinite-dimensional Hilbert spaces and assume that $E \hookrightarrow H \hookrightarrow E'$, with continuous injections, where $E'$ denotes the dual space of $E$. For simplicity, we assume that the continuous injection $E \hookrightarrow H$ has norm at most $1$ and identify $E$ with its image in $H$. Set
$$
\begin{cases} 
\|u\|^2 = \langle u, u \rangle, & u \in E, \\
|u|^2 = (u, u), &  u \in H,
\end{cases}
$$
and for $\mu > 0$, define the set
$$S_\mu = \{u \in E | \, |u|^2 = \mu\}.$$ Clearly, $S_\mu$ is a smooth submanifold of $E$ of codimension $1$, and its tangent space at  $u \in S_\mu$ can be considered as the closed codimension $1$ subspace of $E$ given by
$$T_u S_\mu = \{ v \in E | (u, v) = 0 \}.$$
In the following,  $\| \cdot \|_*$ and $\| \cdot \|_{**}$ denote the operator norms on $L(E, \mathbb{R})$ and  $L(E, L(E, \mathbb{R}))$, respectively.
\begin{definition}[{\cite[Definition 2.1]{Bort}}]\label{defholder}
    Let $\phi : E \to \mathbb{R}$ be a $C^2$-functional on $E$ and $\alpha \in (0, 1]$. We say that $\phi'$ and $\phi''$ are $\alpha$-H\"older continuous on bounded sets if for any $R > 0$, one can find $M = M(R) > 0$ such that, for any $u_1, u_2 \in B(0, R)$,
$$\| \phi'(u_1) - \phi'(u_2) \|_* \leq M \| u_1 - u_2 \|^\alpha, \quad \| \phi''(u_1) - \phi''(u_2) \|_{**} \leq M \| u_1 - u_2 \|^\alpha. $$
\end{definition} 
\begin{definition}[{\cite[Definition 1.3]{Bort0}}]
Let $\phi$ be a $C^2$-functional on $E$. For any $u \in E\backslash\{0\}$,  define the continuous bilinear map\purple{:}\blue{}
$$D^2 \phi(u) = \phi''(u) - \frac{\phi'(u) \cdot u}{\abs{u}^2} (\cdot, \cdot).$$
\end{definition}
\begin{definition}[{\cite[Definition 1.4]{Bort0}}]Let $\phi$ be a $C^2$-functional on $E$. For any $u \in S_\mu$ and $\theta > 0$, we define the {\it approximate Morse index} by
$$\hat{m}_\theta (u) = \sup \left\{ \dim L \bigg| \, L \, \text{is a subspace of } T_u S_\mu \, \text{such that } D^2 \phi(u)[\varphi, \varphi] < -\theta \|\varphi\|^2, \, \forall \varphi \in L \setminus \{0\} \right\}.$$
If $u$ is a critical point for the constrained functional $\phi|_{S_\mu}$ and $\theta = 0$, we call $\hat m_0(u)$ the Morse index of $u$ as a constrained critical point.
\end{definition}
In the context of this paper, we take $E = H^1 (\G)$ and $H = L^2 (\G)$. For $\rho \in [\frac{1}{2},1]$, we denote by $E_\rho(\cdot,\G) : H^1(\G) \to \R$ the energy functional associated with the parameterized equations \eqref{eqpdrho}, defined by
\begin{equation}\label{funcrho}
E_\rho(u,\G) := \frac{1}{2}\norm{u'}^{2}_2-\frac{\rho}{p}\sum_{\vv \in \V_{0}}\abs{u(\vv)}^p.
\end{equation}

Since $p>4$, the maps $E_\rho'$ and $E_\rho''$ are locally Lipschitz on bounded subsets of $H^1(\G)$, uniformly for $\rho\in[1/2,1]$. Hence they are $\alpha$-H\"older continuous on bounded sets in the sense of Definition \ref{defholder}, with $\alpha=1$.

The following two abstract theorems were established in \cite{Bort0}. In the proofs of Theorems \ref{th1} and \ref{th2}, these results are used to obtain bounded Palais--Smale sequences of $E_\rho(\cdot,\G)$ constrained on $H_\mu^1(\G)$ for almost every $\rho \in [\frac{1}{2},1]$.
\begin{theorem}[{\cite[Theorem 1.10]{Bort0}}]\label{minimax0}
Let \( I \subset (0,+\infty) \) be an interval and consider a family of \( C^2 \) functionals \(\Phi_\rho : E \to \mathbb{R} \) of the form
\[
\Phi_\rho(u) = A(u) - \rho B(u), \quad \rho \in I,
\]
where \( B(u) \geq 0 \) for every \( u \in E \), and
\begin{equation}\label{asymptotic0}
    \text{either } A(u) \to +\infty \,\, \text{or } B(u) \to +\infty, \,\, \text{as } \,\|u\| \to +\infty.
\end{equation}
Suppose moreover that \(\Phi_\rho'\) and \(\Phi_\rho''\) are \(\alpha\)-H\"older continuous on bounded sets for some \(\alpha \in (0,1]\). Finally, suppose that there exist \( w_1, w_2 \in S_\mu \) (independent of \( \rho \)) such that, setting
\[
\Gamma = \left\{ \gamma \in C([0,1], S_\mu) \mid \gamma(0) = w_1, \quad \gamma(1) = w_2 \right\},
\]
we have
\begin{equation}\label{strit-ineq0}
    c_\rho := \inf_{\gamma \in \Gamma} \max_{t \in [0,1]} \Phi_\rho(\gamma(t)) > \max \left\{ \Phi_\rho(w_1), \Phi_\rho(w_2) \right\}, \quad \rho \in I.
\end{equation}
Then, for almost every \( \rho \in I \), there exist sequences \(\{u_n\} \subset S_\mu \) and \(\zeta_n \to 0^+\) such that, as \( n \to +\infty \),
\begin{enumerate}[label=\rm(\roman*)]
\item \(\Phi_\rho(u_n) \to c_\rho;\)
\item \(\|\Phi_\rho'|_{S_\mu}(u_n)\|_* \to 0;\)\label{minimax02}
\item \(\{u_n\}\) is bounded in \(E;\)
\item \(\hat{m}_{\zeta_n}(u_n) \leq 1.\)
\end{enumerate}
\end{theorem}
\begin{theorem}[{\cite[Theorem 1.12]{Bort0}}]\label{minimax}
	Let \( I \subset (0, \infty) \) be an interval and consider a family of \( C^2 \) functionals \( \Phi_{\rho} : E \to \mathbb{R} \) of the form
	\[
	\Phi_\rho(u) = A(u) - \rho B(u), \quad \forall \rho \in I,
	\]
	where \( B(u) \geq 0 \) for all \( u \in E \) and
    \begin{align}\label{asymptotic}
  \text{either }   A(u) \to +\infty \text{ or } B(u) \to +\infty, \quad \text{as }  \|u\| \to +\infty.
    \end{align}
	Suppose that, for every \( \rho \in I \), \( \Phi_\rho|_{S_\mu} \) is even and moreover that \( \Phi'_\rho \) and \( \Phi''_\rho \) are \( \alpha \)-H\"older continuous on bounded sets for some \( \alpha \in (0, 1] \). Finally, suppose that there exists an integer \( N \geq 2 \) and two odd maps \(\gamma_i:\mathbb{S}^{N-2}\to S_\mu\), \(i=0,1\), such that the set\begin{align}\label{GamN}
	\Gamma_N := \left\{ \gamma \in C([0, 1] \times \mathbb{S}^{N-2}, S_\mu) \mid \forall t \in [0, 1], \gamma(t, \cdot) \text{ is odd, } \gamma(0, \cdot) = \gamma_0 \text{ and } \gamma(1, \cdot) = \gamma_1 \right\}
	\end{align}
	is nonempty and that
\begin{align}\label{strit-ineq}
c_\rho^N := \inf_{\gamma \in \Gamma_N} \max_{(t,a) \in [0, 1] \times \mathbb{S}^{N-2}} \Phi_\rho(\gamma(t, a)) > \max_{a \in \mathbb{S}^{N-2}} \left\{ \Phi_\rho(\gamma_0(a)), \Phi_\rho(\gamma_1(a)) \right\}, \quad \forall \rho \in I.
\end{align}
	Then, for almost every \( \rho \in I \), there exist sequences \(\{u_n\} \subset S_\mu\) and \(\zeta_n \to 0^+\) such that, as \( n \to +\infty \),
	\begin{enumerate}[label=\rm(\roman*)]
		\item \( \Phi_\rho(u_n) \to c_\rho^N \);
		\item \( \|\Phi'_\rho|_{S_\mu}(u_n)\|_* \to 0 \);\label{minimax2}
		\item \(\{u_n\}\) is bounded in \( E \);
		\item \( \hat{m}_{\zeta_n}(u_n) \leq N \).
	\end{enumerate}
\end{theorem}
By \cite[Remark 1.6]{Bort0}, Theorem \ref{minimax0} \ref{minimax02} (or Theorem \ref{minimax} \ref{minimax2}) implies the existence of a sequence $\{\lambda_n\} \subset \mathbb{R}$ such that
\begin{equation}
\Phi^{ \prime}_{\rho}(u_n, \mathcal{G}) + \lambda_n (u_n, \cdot) \to 0 \quad \text{in the dual of $H^1(\mathcal{G})$ },
\end{equation}
where the sequence of Lagrange multipliers $\{\lambda_n\}$ is defined by
\begin{equation}
\lambda_n = -\frac{1}{\mu} \Phi'_\rho(u_n, \mathcal{G})[u_n].
\end{equation}
To apply Theorem \ref{minimax0} (or Theorem \ref{minimax}), we consider $I=[\frac{1}{2},1]$, $A(u)=\frac{1}{2}\norm{u'}^{2}_2$ and $B(u)=\frac{1}{p}\sum_{\vv \in \V_{0}}\abs{u(\vv)}^p$ and $\Phi_\rho(u)=E_\rho(u,\G)$. Since
$$
u \in H^1_\mu(\G), \quad \norm{u}\to +\infty \Rightarrow A(u) \to +\infty,
$$
it follows that assumptions \eqref{asymptotic0} and \eqref{asymptotic} are satisfied.





\iffalse
At the end of this subsection, we recall the following result from \cite{Ca}, which plays a vital role in proving that the Lagrange multipliers sequences of the bounded Palais-Smale sequences, obtained by the above theorems, converge to nonnegative numbers. 
\begin{lemma}[{\cite[Lemma 2.7]{Ca}}]\label{lem 2.7}
Let $\Phi_\rho$ be defined as in Theorem \ref{minimax0} (or Theorem \ref{minimax}), \(\{u_{n}\}\subset S_{\mu}\), \(\{\lambda_{n}\}\subset\mathbb{R}\) and \(\{\zeta_{\nu}\}\subset\mathbb{R}^{+}\) with \(\zeta_{\nu}\to 0^{+}\). Assume that, for a given \(M\in\mathbb{N}^+\), the following conditions hold:
\begin{enumerate}[label=\rm(\roman*)]
\item For large enough \(n\in\mathbb{N}^+\), all subspaces \(W_{n}\subset E\) with the property
\[
\Phi_\rho^{\prime\prime}(u_{n})[\varphi,\varphi]+\lambda_{n}|\varphi|^{2 }<-\zeta_{n}\|\varphi\|^{2},\qquad\text{for all }\varphi\in W_{n}\setminus\{0\},
\]
satisfy: \(\dim W_{n}\leq M\).
\item There exist \(\lambda\in\mathbb{R}\), a subspace \(Y\) of \(E\) with \(\dim Y\geq M+1\) and \(\zeta>0\) such that, for large enough \(n\in\mathbb{N}^+\),
\begin{equation}\label{eq2.8}
\Phi_\rho^{\prime\prime}(u_{n})[\varphi,\varphi]+\lambda|\varphi|^{2} \leq-\zeta\bigl{\|}\varphi\bigr{\|}^{2},\qquad\text{for all }\varphi\in Y.
\end{equation}
\end{enumerate}
Then \(\lambda_{n}>\lambda\) for all large enough \(n\in\mathbb{N}^+\). In particular, if \eqref{eq2.8} holds for any \(\lambda<0\), then \(\liminf_{n\to\infty}\lambda_{n}\geq 0\).
\end{lemma}

\fi

\subsection{Boundedness of the sequence of solutions to the parameterized equations}
In this subsection, we establish a key compactness lemma for passing from the parameterized problems \eqref{eqpdrho} to the original problem \eqref{eqpd}. More precisely, if $\rho_n\to1^-$ and $\{(u_{\rho_n},\lambda_{\rho_n})\}\subset H^1(\G)\times(0,+\infty)$ is a sequence of normalized solutions of \eqref{eqpdrho} with bounded energy, then $\{u_{\rho_n}\}$ is a bounded Palais--Smale sequence for $E(\cdot,\G)$ constrained on $H^1_\mu(\G)$.
\begin{lemma}\label{lembu}
Let $\G$ be a noncompact metric graph with finitely many edges, let $p > 4$, and let $ \{\rho_n\} \subset [\frac{1}{2}, 1] $ be a sequence such that $\rho_n\rightarrow 1^{-}$. Suppose that 
$$\{(u_{\rho_n},\lambda_{\rho_n})\} \subset H^1(\G)\times (0,+\infty)$$
is a sequence of solutions  to 
\begin{equation}\label{eqpdn}
    \begin{cases} 
u_{\rho_n}'' = \lambda_{\rho_n} u_{\rho_n} & \text{on every }\e \in \E, \\ 
\int_\G\abs{u_{\rho_n}}^2\, dx = \mu & \\
\displaystyle\sum_{\e \succ \vv} u'_{\rho_n,\e}(\vv) = -\rho_n|u_{\rho_n}(\vv)|^{p-2}u_{\rho_n}(\vv) & \text{at every } \vv \in \V_0,\\
 \displaystyle\sum_{\e \succ \vv} u'_{\rho_n,\e}(\vv) = 0 & \text{at every } \vv \in \V \setminus \V_0,
\end{cases}
\end{equation}
where $u_{\rho_n,\e}'(\vv)$ stands for $u_{\rho_n}|_\e'(0)$ or $-u_{\rho_n}|_\e'(\ell_\e)$, according to whether the vertex $\vv$ is identified with $0$ or $\ell_\e$. Assume additionally  that the sequence of energy levels $\{c_{\rho_n} := E_{{\rho_n}}(u_{\rho_n}, \G)\}$ is bounded. Then, $\{\lambda_{\rho_n}\} \subset (0,+\infty)$ is bounded from above. Moreover, up to a subsequence, $\{u_{\rho_n}\}$ satisfies the following properties.
 %is a bounded Palais-Smale sequence for $E(\cdot, \G)$ constrained on $H_\mu^1(\G)$ satisfying 
	\begin{enumerate}[label=\rm(\roman*)]
		\item $E(u_{\rho_n},\G) \to c$, where  $\displaystyle c:= \lim_{n \to \infty}c_{\rho_n}$;\label{bps1}
		\item $E'(u_{\rho_n},\G)+\lambda_{\rho_n}(u_{\rho_n},\cdot)_2 \to 0 $ \text{in the dual of $H^1(\mathcal{G})$};
		\item $\{u_{\rho_n}\}$ is bounded in $ H^1(\G) $.\label{bps3}
	\end{enumerate}
\end{lemma}
\begin{proof}
Let $\{(u_{\rho_n},\lambda_{\rho_n})\} \subset H^1(\G)\times (0,+\infty)$ be a sequence of solutions to \eqref{eqpdn}, where $ (\rho_n)$ is a sequence in $[\frac{1}{2},1]$ satisfying $\rho_n \to 1^{-}$ as $n \to \infty$. We first show that  $\{\lambda_{\rho_n}\}$ cannot diverge to $+\infty$. Assume by contradiction that $\lambda_{\rho_n} \to +\infty$ as $n \to \infty$. By a straightforward computation, we have
$$
\begin{aligned}
    2c_{\rho_n}&=2E_{\rho_n}(u_{\rho_n},\G)-E_{\rho_n}'(u_{\rho_n} ,\G)[u_{\rho_n}]-\lambda_{\rho_n}(u_{\rho_n},u_{\rho_n})_2\\
    &=\frac{\rho_n(p-2)}{p}\sum_{\vv \in \V_0}\abs{u_{\rho_n}(\vv)}^p-\lambda_{\rho_n}\mu\\
    &\geq \frac{p-2}{2p}\sum_{\vv \in \V_0}\abs{u_{\rho_n}(\vv)}^p-\lambda_{\rho_n}\mu,
\end{aligned}
$$
and thus, there exists a constant $C_1>0$  such that, for all sufficiently large $n\in\mathbb{N}^+$, 
\begin{equation}\label{eqc1lambdan0}
    \abs{u_{\rho_n}(\vv)} \leq \left(\sum_{\vv \in \V_0}\abs{u_{\rho_n}(\vv)}^p\right)^\frac{1}{p} \leq C_1\lambda_{\rho_n}^\frac{1}{p}, \quad \forall \vv \in \V_0.
\end{equation}
Since \(u_{\rho_n}\in H^1(\G)\), its restriction to each half-line of \(\G\)  tends to zero at infinity. Moreover, as \(\G\) has only finitely many edges, it follows that \(|u_{\rho_n}|\) attains its maximum at some point \(x_n\in\G\), hence
$|u_{\rho_n}(x_n)|=\norm{u_{\rho_n}}_\infty$.

Now, we prove that $x_n \in \V_0$. Without loss of generality, we may assume that $u_{\rho_n}(x_n)>0$.  If $x_n$ is in the interior of some edge $e \in \E$, then we have $u''_{\rho_n}(x_n)/u_{\rho_n}(x_n) \leq 0$, which is a contradiction since $u''_{\rho_n}(x_n) = \lambda_{\rho_n}u_{\rho_n}(x_n)$ by \eqref{eqpdn}. Thus, we obtain $x_n \in \V$. We argue by contradiction and suppose that $x_n \in \V \setminus \V_0$. After parametrizing every edge incident at $x_n$ so that $x_n$ corresponds to $0$, we have $u'_{\rho_n,\e}(x_n)\leq0$ for every $\e\succ x_n$. Thus, since $\displaystyle\sum_{\e \succ x_n} u'_{\rho_n,\e}(x_n)=0$ by \eqref{eqpdn}, we obtain $u'_{\rho_n,\e}(x_n)=0$ for every $\e \succ x_n$. Therefore, since $u''_{\rho_n}(x_n)=\lambda_{\rho_n}u_{\rho_n}(x_n)>0$, we have, for every $\e \succ x_n$, $u'_{\rho_n,\e}(x)>0$ for all $x \in \e$ such that $\operatorname{dist}(x,x_n)>0$ small enough, where $x_n$ is identified with $0$. This implies that $u_{\rho_n,\e}(x) > u_{\rho_n,\e}(x_n)$ for some $x \in \e$ and $\e \succ x_n$, which is a contradiction. Hence, we conclude that  $x_n \in \V_0$. Thus, by \eqref{eqc1lambdan0},
     \begin{equation}\label{eqc1lambdan}
                 \abs{u_{\rho_n}(\vv)} \leq \abs{u_{\rho_n}(x_n)} \leq C_1\lambda_{\rho_n}^\frac{1}{p}, \quad \forall \vv \in \V.
     \end{equation}
    
On one hand, for any unbounded edge $\e \in \E$, it follows from \cite[Remark 4 in Section 3.3]{AD} and $\int_\G\abs{u_{\rho_n}}^2\, dx = \mu$ that
$$
u_{\rho_n,\e}(x)=u_{\rho_n,\e}(0)e^{-\sqrt{\lambda_{\rho_n}}x},
$$
where $u_{\rho_n,\e}:=u_{\rho_n}|_\e$. Thus, by \eqref{eqc1lambdan} and $p >4$, we conclude
$$
\int_{I_\e} \abs{u_{\rho_n,\e}}^2 \, dx =\frac{\abs{u_{\rho_n,\e}(0)}^2}{2\sqrt{\lambda_{\rho_n}}} \leq \frac{C_1^2}{2}\lambda_{\rho_n}^{\frac{2}{p}-\frac{1}{2}}\to 0 \quad \text{as }n \to +\infty.
$$
On the other hand, for any $n \in \mathbb{N}^+$ and any bounded edge $\e \in \E$, it follows from \cite[Remark 4 in Section 3.3]{AD} that there exist two constants $a_{\rho_n,\e}, b_{\rho_n,\e} \in \R$ such that  
$$
u_{\rho_n,\e}(x)=a_{\rho_n,\e} e^{\sqrt{\lambda_{\rho_n}}x} + b_{\rho_n,\e} e^{-\sqrt{\lambda_{\rho_n}}x}.
$$
Then,
$$
a_{\rho_n,\e} = \frac{u_{\rho_n,\e}(\ell_\e)e^{\sqrt{\lambda_{\rho_n}}\ell_\e}-u_{\rho_n,\e}(0)}{e^{2\sqrt{\lambda_{\rho_n}}\ell_\e}-1}\quad  \text{and} \quad b_{\rho_n,\e} =\frac{u_{\rho_n,\e}(0)-u_{\rho_n,\e}(\ell_\e)e^{-\sqrt{\lambda_{\rho_n}}\ell_\e}}{1-e^{-2\sqrt{\lambda_{\rho_n}}\ell_\e}},
$$
and thus, by \eqref{eqc1lambdan}, we know that there exists a constant $C_2>0$ such that, for all sufficiently large $n\in\mathbb{N}^+$,
$$
\abs{a_{\rho_n,\e}} \leq C_2 \lambda_{\rho_n}^\frac{1}{p}e^{-\sqrt{\lambda_{\rho_n}}\ell_\e}\quad  \text{and} \quad \abs{b_{\rho_n,\e}} \leq C_2 \lambda_{\rho_n}^\frac{1}{p},$$
whence it follows from $p>4$ that
$$
\int_{I_\e} \abs{u_{\rho_n,\e}}^2 \, dx=2a_{\rho_n,\e} b_{\rho_n,\e} \ell_\e +\frac{a_{\rho_n,\e}^2(e^{2\sqrt{\lambda_{\rho_n}}\ell_\e} -1)}{2\sqrt{\lambda_{\rho_n}}} - \frac{b_{\rho_n,\e}^2(e^{-2\sqrt{\lambda_{\rho_n}}\ell_\e} -1)}{2\sqrt{\lambda_{\rho_n}}} \to 0, \quad \text{as }n \to \infty.
$$

Now, since $\e \in \E$ is arbitrary and $\E$ is a finite set, we have $\int_\G\abs{u_{\rho_n}}^2\, dx \to 0$ as $n \to \infty$, contradicting the fact that $\int_\G\abs{u_{\rho_n}}^2\, dx =\mu $ for all $n \in \mathbb{N}^+$. Therefore, the sequence $\{\lambda_{\rho_n}\}$ is bounded from above.

Next, we compute that 
$$
pc_{\rho_n}=pE_{\rho_n}(u_{\rho_n},\G)-E_{\rho_n}'(u_{\rho_n} ,\G)[u_{\rho_n}]-\lambda_{\rho_n}(u_{\rho_n},u_{\rho_n})_2=(\frac{p}{2}-1)\norm{u_{\rho_n}'}_2^2-\lambda_{\rho_n}\mu,
$$
which shows that $\{u_{\rho_n}\}$ is bounded in $H^1(\G)$. Thus,
$$
E(u_{\rho_n},\G)-E_{\rho_n}(u_{\rho_n},\G)\to 0  \quad \text{and}\quad E'(u_{\rho_n},\G)-E'_{\rho_n}(u_{\rho_n},\G)\to 0 \quad \text{as }n \to \infty.
$$
Therefore, up to a subsequence, $\{u_{\rho_n}\}$ is a bounded Palais-Smale sequence for $E(\cdot, \G)$ constrained on $H_\mu^1(\G)$, satisfying properties \ref{bps1}--\ref{bps3}.
\end{proof}


\section{Compactness of bounded Palais-Smale sequences}\label{sec3}
In this section, we assume that $\rho \in [\frac{1}{2},1]$. Let $\{u_n\} \subset H_\mu^1(\mathcal{G})$ be a bounded Palais-Smale sequence for $E_\rho(\cdot, \G)$ constrained on $H_\mu^1(\G)$ at some level $c>0$, that is, $\{u_n\}$ is bounded in $H^1(\mathcal{G})$ and satisfies
\begin{equation}\label{PS1}
E_\rho(u_n, \mathcal{G}) \to c 
\end{equation}
and 
\begin{equation}\label{PS2}
E^{ \prime}_{\rho}(u_n, \mathcal{G}) + \lambda_n (u_n, \cdot)_2 \to 0 \quad \text{in the dual of $H^1(\mathcal{G})$},
\end{equation}
where $\{\lambda_n\} \subset \R$ is the sequence of Lagrange multipliers.


We now prove the compactness of the bounded Palais-Smale sequences $\{u_n\} \subset H^1(\mathcal{G})$.
\begin{lemma}\label{lemurhosolve}
There exists \( u_\rho \in H^1(\mathcal{G}) \) such that, up to a subsequence,
\[
u_n \rightharpoonup  u_\rho \quad \text{in } H^1(\mathcal{G})
\]
and
\[
u_n(x) \to u_\rho(x), \quad \forall x \in \G. 
\]
Moreover, $u_\rho$ is a solution of equation
\begin{equation}\label{eqpdnomu}
\begin{cases} 
u_{\rho}''=\lambda_\rho u_\rho &\quad \text{on every } \e\in \E,\\
\sum_{\e\succ\vv}u_{\rho,\e}'(\vv)=-\rho|u_{\rho}(\vv)|^{p-2}u_{\rho}(\vv)  &\quad \text{at every } \vv\in\V_0,\\
\sum_{\e\succ\vv}u_{\rho,\e}'(\vv)=0  &\quad \text{at every } \vv\in \V \setminus \V_0,
\end{cases}
\end{equation}
where $\displaystyle \lambda_\rho=\lim_{n \to +\infty} \lambda_n $ and $u_{\rho,\e}'(\vv)$ stands for $u_\rho|_\e'(0)$ or $-u_\rho|_\e'(\ell_\e)$, according to whether the vertex $\vv$ is identified with $0$ or $\ell_\e$.
\end{lemma}
\begin{proof}
Since \(\{u_n\} \subset H^1(\mathcal{G})\) is bounded, by \cite[Theorem 8.8]{Ha},  we may assume that, up to a subsequence, there exists \( u_\rho \in H^1(\mathcal{G}) \) such that
\[
u_n \rightharpoonup  u_\rho \quad \text{in } H^1(\mathcal{G})
\]
and
%\begin{equation*}
%u_n \to u_\rho \quad \text{in } L^r_{\text{loc}}(\mathcal{G}) \quad \text{for all } r \geq 2,
%\end{equation*}
\[
u_n(x) \to u_\rho(x), \quad \forall x \in \G.
\]
Observe also that, from \eqref{PS2} and the boundedness of \(\{u_n\} \subset H^1(\mathcal{G})\), we obtain
$$
   o(1)= E_{\rho}'(u_n ,\G)[u_n]+\lambda_n(u_n,u_n)_2 =E_{\rho}'(u_n ,\G)[u_n]+\lambda_n \mu,
$$
which implies that $\{\lambda_n\}$ is bounded. Thus, up to a subsequence, there exists \(\lambda_\rho \in \mathbb{R}\) such that 
\[
\lambda_n \rightarrow \lambda_\rho,\,\,\,\text{as}\,\,\,n \to +\infty.
\]
Next, we show that $u_{\rho}\in H^{1}(\G)$ is a weak solution of \eqref{eqpdnomu}. Indeed, using \eqref{PS2} and  $\displaystyle \lim_{n \to +\infty} \lambda_n = \lambda_\rho$, we have, for every \(\eta\in H^{1}(\mathcal{G})\),
\begin{equation}\label{solution u}
\begin{split}
0&=\lim_{n\to\infty}\left(E^{ \prime}_{\rho}(u_{n},\mathcal{G})+\lambda_{n}(u_{n},\cdot)_2\right)[\eta]\\
&=\lim_{n\to\infty}\left(\int_{\mathcal{G}}u_{n}^{\prime}\eta^{ \prime}\,dx+\lambda_{n}\int_{\mathcal{G}}u_{n}\eta\,dx-\rho\sum_{\vv\in \V_0} |u_{n}(\vv)|^{p-2 }u_{n}(\vv)\eta(\vv)\right)\\
&=\int_{\mathcal{G}}(u_{\rho})^{\prime}\eta^{\prime}\,dx+\lambda_{\rho}\int_{\mathcal{G}}u_{\rho}\eta\,dx-\rho\sum_{\vv\in \V_0}|u_{\rho}(\vv)|^{p-2}u_{\rho}(\vv)\eta(\vv).
\end{split}
\end{equation}
This completes the proof of the lemma.
\end{proof}
Next, we focus on proving that $u_{n} \to u_\rho$ in $H^1(\G)$, as $n \to \infty$, to ensure that the limit \(u_\rho\) satisfies the mass constraint, that is, \(u_\rho\in H^1_\mu(\mathcal{G})\).

\begin{lemma}\label{lemurhoconverge}
 Assume that $\lambda_\rho \geq 0$ in  Lemma \ref{lemurhosolve}. Then,
\begin{equation}\label{strong conv}
\int_{\mathcal{G}}|(u_{n}-u_{\rho})^{\prime}|^{2}\,dx+\lambda_{\rho}\int_{\mathcal{G}}|u_{n}-u_{\rho}|^{2}\,dx\to 0 \quad \text{as }n \to \infty.
\end{equation}
In particular, if \(\lambda_{\rho}>0\), the sequence \(\{u_{n}\}\) converges strongly to $u_{\rho}$ in \(H^{1}(\mathcal{G})\).
\end{lemma}

\begin{proof}
Using \eqref{PS2} and the fact that $\displaystyle\lim_{n \to +\infty} \lambda_n = \lambda_\rho$, for every \(\eta\in H^{1}(\mathcal{G})\), we have
\begin{align*}
o(1)\|\eta\|&=\left(E^{ \prime}_{\rho}(u_{n},\mathcal{G})+\lambda_{n}(u_{n},\cdot)_2\right)[\eta]\\
&= \int_{\mathcal{G}}u_{n}^{\prime}\eta^{\prime}\,dx-\rho\sum_{\vv\in \V_0}|u_n(\vv)|^{p-2}u_n(\vv)\eta(\vv)+\lambda_{n}\int_{\mathcal{G}}u_{n}\eta\,dx\\
&=\int_{\mathcal{G}}u_{n}^{\prime}\eta^{\prime}\,dx-\rho\sum_{\vv\in \V_0}|u_{\rho}(\vv)|^{p-2}u_{\rho}(\vv)\eta(\vv)+\lambda_{\rho}\int_{\mathcal{G}}u_{n}\eta\,dx+(\lambda_{n}-\lambda_{\rho})\int_{\mathcal{G}}u_{n}\eta\,dx+o(1)\|\eta\|,
\end{align*}
that is,
\begin{equation}\label{PS21}
\int_{\mathcal{G}}u_{n}^{\prime}\eta^{\prime}\,dx-\rho\sum_{\vv\in \V_0}|u_{\rho}(\vv)|^{p-2}u_{\rho}(\vv)\eta(\vv)+\lambda_{\rho}\int_{\mathcal{G}}u_{n}\eta\,dx=o(1)\|\eta\|.
\end{equation}

Now, taking the difference between \eqref{PS21} and \eqref{solution u}, choosing \(\eta=\eta_{n}:=u_{n}-u_{\rho}\), we obtain
\begin{align*}
o(1)=o(1)\|u_{n}-u_{\rho}\|
&=\int_{\mathcal{G}}(u_{n}^{\prime}-u_{\rho}^{\prime})\eta_{n}^{\prime}\,dx+\lambda_{\rho}\int_{\mathcal{G}}(u_{n}-u_{\rho})\eta_{n}\,dx\\
&=\int_{\mathcal{G}}\big{|}(u_{n}-u_{\rho})^{\prime}\big{|}^{2 }\,dx+\lambda_{\rho}\int_{\mathcal{G}}|u_{n}-u_{\rho}|^{2}\,dx,
\end{align*}
which yields \eqref{strong conv}.
\end{proof}

\begin{lemma}\label{lem 5.2}
For any \(\lambda<0\) and \(d\in\mathbb{N}^+\), there exists a subspace \(Y\subset H^{1}(\mathcal{G})\) with \(\dim Y =d\) such that
\[
%\int_\mathcal{G} |w'|^2 dx+ \lambda\int_\mathcal{G} w^2  dx-(p-1)\rho\sum _{\vv\in \V}|u_n(\vv)|^{p-2}w^{2}(\vv)= 
E''_\rho(u_n, \mathcal{G})[w, w] + \lambda \|w\|_2^2 \leq\frac{\lambda}{2}\|w\|^{2}, \qquad\forall\,w\in Y.
\]
\end{lemma}
\begin{proof}
Choose $\varphi\in C_c^\infty(\mathbb{R})$ such that $\operatorname{supp}\varphi\subset(0,1)$ and $\int_0^{+\infty}|\varphi|^2\,dx=1$. For \(t\in\mathbb{R}^{+}\), define the function \(\varphi^{t}\) by
\begin{align}\label{scal}
	\varphi^{t}(x):=t^{1/2}\varphi(tx).
\end{align}
We view $\varphi$ as a function in $H^1(\G)$ whose support is contained in a half-line identified with $[0,\infty)$ and define
\[
\varphi_{1}:=\varphi^{\tau},
\]
where $\tau>0$ is chosen small enough so that
\begin{equation}\label{5.9}
\tau^{2}\|\varphi^{\prime}\|_{L^{2}(\mathbb{R})}^{2}+\lambda\leq\frac{\lambda}{2}(\tau^{2}\|\varphi^{\prime}\|_{L^{2}(\mathbb{R})}^{2}+1).
\end{equation}
We have
\[
\|\varphi_{1}\|_{L^{2}({\G})}^{2}=1,\qquad\|\varphi^{\prime}_{1}\|_{L^{2}({\G})}^{2}=\tau^{2}\|\varphi^{\prime}\|_{L^{2}({\G})}^{2}.
\]
For $i\geq 2$, define
\[
\varphi_{i}(x):=\varphi_{1}\left(x-\frac{i-1}{\tau}\right).
\]
Since $\operatorname{supp}(\varphi_{i})\subset\left(\frac{i-1}{\tau},\frac{i}{\tau}\right)$, the functions $\varphi_{i}$ have mutually disjoint supports. Let $Y=\text{span}\{\varphi_{1},\ldots,\varphi_{d}\}$. Every $w\in Y$ can be written as
\[
w=\sum_{i=1}^{d}\theta_{i}\varphi_{i},\,\,\,\theta_{i}\in \mathbb{R}.
\]
By a direct calculation, we have
\begin{align*}
\int_{\G}|w^{\prime}|^{2}\,dx+\lambda\int_{\G}|w|^{2}\,dx 
&= \tau^{2}\Bigg{(}\sum_{i=1}^{d}\theta_{i}^{2}\|\varphi^{\prime}\|_{L^{2}(\mathbb{R})}^{2}\Bigg{)}+\lambda\left(\sum_{i=1}^{d}\theta_{i}^{2}\right)\\
&=(\tau^{2}\|\varphi^{\prime}\|_{L^{2}(\mathbb{R})}^{2}+\lambda)\sum_{i=1}^{d}\theta_{i}^{2},
\end{align*}
and similarly, 
$$\|w\|^{2}=(\tau^{2}\|\varphi^{\prime}\|_{L^{2}(\mathbb{R})}^{2}+1)\sum_{i=1}^{d}\theta_{i}^{2}.$$
Therefore, \eqref{5.9} and $\operatorname{supp}(\varphi_{i})\subset\left(\frac{i-1}{\tau},\frac{i}{\tau}\right)$ imply that
\begin{align*}
\int_\mathcal{G} |w'|^2 \,dx+ \lambda\int_\mathcal{G} w^2 \,dx-(p-1)\rho\sum _{\vv\in \V_0}|u_n(\vv)|^{p-2}w^{2}(\vv)
&= (\tau^{2}\|\varphi^{\prime}\|_{L^{2}(\mathbb{R})}^{2}+\lambda)\sum_{i=1}^{d}\theta_{i}^{2}\\
&\leq \frac{\lambda}{2}(\tau^{2}\|\varphi^{\prime}\|_{L^{2}(\mathbb{R})}^{2}+1)\sum_{i=1}^{d}\theta_{i}^{2}\\
&= \frac{\lambda}{2}\|w\|^{2}.
\end{align*}
This completes the proof of Lemma \ref{lem 5.2}.
\end{proof}

\begin{lemma}\label{lamuna+}
Assume that there exists a sequence $\{\zeta_n\} \subset \mathbb{R}^+$ with $\zeta_n \to 0^+$ such that  $ \hat{\mathrm{m}}_{\zeta_n}(u_n) \leq N $, that is, if
\begin{equation}\label{5.4}
\int_\mathcal{G} |\varphi'|^2 \,dx+ \lambda_n\int_\mathcal{G} \varphi^2  \,dx-(p-1)\rho\sum _{\vv\in \V_0}|u_n(\vv)|^{p-2}\varphi^{2}(\vv)= E''_\rho(u_n, \mathcal{G})[\varphi, \varphi] + \lambda_n \|\varphi\|_2^2 < -\zeta_n \|\varphi\|^2 
\end{equation}
 holds for every \(\varphi\in W_n\setminus\{0\}\), where \(W_n\subset T_{u_n}H^1_\mu(\mathcal{G})\), \blue{then } one has \(\dim W_n\leq N\). Then $\lambda_\rho\geq0$.
\end{lemma}
\begin{proof}
    Observe that the codimension of $T_{u_{n}}H^{1}_{\mu}({\G})$ in $H^{1}({\G})$ is one. Therefore, if inequality \eqref{5.4} holds for all $\varphi\in W_{n}\setminus\{0\}$ in a subspace $W_{n}\subset H^{1}({\G})$, then  \(\dim W_n \leq  N+1 \). Let $\lambda<0$ and $Y$ be the subspace of dimension $d=N+2$ provided by Lemma \ref{lem 5.2}. Applying \cite[Lemma 2.7]{Ca}, we conclude that
\begin{equation}\label{5.10}
\lambda_{\rho}\geq 0.
\end{equation}
\end{proof}


\section{Proof of Theorem \ref{th1}}\label{sect4}
In this section, we prove Theorem \ref{th1} by applying the abstract minimax principle stated in Theorem \ref{minimax0}. To ensure that the family $\{E_{\rho}(\cdot,\G)\}$ fits into the framework of Theorem \ref{minimax0}, we first show that it has a mountain pass geometry on $H^{1}_{\mu}(\G)$ uniformly with respect to $\rho\in [\frac{1}{2},1]$.
\begin{lemma}\label{lem3.1}
For every \(\mu>0\), there exist \(w_{1},w_{2}\in H_{\mu}^{1}(\G)\) with $w_1,w_2\geq 0$ independent of \(\rho\in\left[\frac{1}{2},1\right]\) such that  
\[
c_{\rho}:=\inf_{\gamma\in\Gamma}\max_{t\in[0,1]}E_{\rho}(\gamma(t), \mathcal{G})>\max\{E_{\rho}(w_{1},\mathcal{G}),E_{\rho}(w_{2}, \mathcal{G})\},\qquad\forall\rho\in\left[\frac{1}{2},1\right],
\]
where \(\Gamma:=\left\{\gamma\in C([0,1],H^{1}_{\mu}(\mathcal{G}))\,\mid \gamma(0)=w_{1},\gamma(1)=w_{2}\right\}\).
\end{lemma}
\begin{proof}
For any \(\mu,k>0\), denote
\begin{equation*}
A_{\mu,k}:=\{u\in H^{1}_{\mu}(\mathcal{G})\,\big{|}\,\int_{\mathcal{G}}|u^{\prime}|^{2}\,dx<k\}, \,\,\,  \partial A_{\mu,k}:=\{u\in H^{1}_{\mu}(\mathcal{G})\,\big{|}\,\int_{\mathcal{G}}|u^{\prime}|^{2}\,dx=k\}.
\end{equation*}
Since $\mathcal{G}$ is noncompact, it contains at least one half-line. Choose $v\in C_c^\infty((0,\infty))$ with $v\geq0$ and $\|v\|_2^2=\mu$, regard it as a function supported on that half-line, and define $v^t(x)=t^{1/2}v(tx)$ for $t>0$. Then $\|v^t\|_2^2=\mu$ and $\|(v^t)'\|_2^2=t^2\|v'\|_2^2$. Hence $v^t\in A_{\mu,k}$ for all sufficiently small $t>0$, while the choice $t=\sqrt{k}/\|v'\|_2$ gives $v^t\in\partial A_{\mu,k}$. Therefore, both $A_{\mu,k}$ and $\partial A_{\mu,k}$ are nonempty.


Now, by the Gagliardo-Nirenberg inequality \eqref{eqgn}, for any $u\in H^{1}_{\mu}(\G)$, we obtain
\[
E_{\rho}(u,\G)\geq\frac{1}{2}\|u^{\prime}\|_2^2-\rho\frac{C_{p,\G} }{p}\mu^{\frac{p}{4}}\norm{u'}_2^{\frac{p}{2}}.
\]
Then, for any \(\mu>0\) and \(u\in\partial A_{\mu,k_{0}}\) with \(k_{0}:=\frac{1}{2}(\frac{p}{2C_{p,\G}})^{\frac{4}{p-4}}\mu^{-\frac{p}{p-4}}\), we have
\[
\inf_{u\in\partial A_{\mu,k_{0}}}E_{\rho}(u,\mathcal{G})\geq k_{0}\left(\frac{1}{2}-\frac{C_{p,\G}}{p}\mu^{\frac{p}{4}}k_{0}^{\frac{p-4}{4}}\right)=\colon\alpha>0,
\]
for every \(\rho\in[\frac{1}{2},1]\). Next, observe that for any \(u\in A_{\mu,k}\)
\[
E_{\rho}(u,\mathcal{G})\leq\frac{1}{2}\int_{\mathcal{G}}|u^{\prime}|^{2}\,dx\leq\frac{1}{2}k.
\]
Thus, recalling that \(A_{\mu,k}\neq\emptyset\) for all \(\mu,k>0\), it is possible to choose a function \(w_{1}\in A_{\mu,k}\)  with $w_1\geq 0$ for $k>0$ sufficiently small, such that
\begin{equation}\label{eq3.3}
\|w_{1}^{\prime}\|_2^2<k_{0}\quad\text{and}\quad E_{\rho}(w_{1},\mathcal{G})<\frac{\alpha}{2}\quad\forall\rho\in\left[\frac{1}{2},1\right].
\end{equation}

Next, take a vertex \(\vv \in \V_0\), that we identify in the following with $0$, and denote by \(\e_j\), \(j = 1, \ldots, K\), the edges incident at \(\vv\). Set
\begin{equation*}
\ell :=\begin{cases}
    \min_{1 \leq j \leq K} |\e_j|&\text{if }\e_j \text{ is bounded for some }j = 1, ..., K,\\
    1\quad &\text{ otherwise}.
\end{cases} 
\end{equation*}
Define \(u_{s_{0}}\in H^1(\mathcal{G})\) by
$$
	u_{s_0}(x) = 
	\begin{cases} 
		s_0(\frac{\ell}{K} - x) & \text{on }  \e_j \cap [0,\frac{\ell}{K}], \quad \forall \, j = 1, \ldots, K,\\
		0 & \text{elsewhere,}
	\end{cases}
$$
with $s_0>0$ chosen so that \(\|u_{s_{0}}\|_2^{2}=\mu\), and denote  $w:=u_{s_{0}}$. For \(t>1\), set \(w^t(x):=t^{1/2}w(tx)\). Then \(w^t\in H^1_\mu(\mathcal{G})\), and
\[
E_{\rho}(w^{t},\mathcal{G})= \frac{t^2}{2}\norm{w'}^2_2-\frac{\rho}{p} t^{p/2} \left ( \frac{s_0\ell}{K}  \right )^{p},
\]
for every \(\rho\in[\frac{1}{2},1]\). Since \(p>4\), $E_{\rho}(w^{t},\mathcal{G})\rightarrow -\infty$ as \(t\to+\infty\). Hence,  there exists a sufficiently large \(t_2>0\) such that,  for all \(t>t_{2}\) 
$$t^{2}\|w^{\prime}\|_2^2>2k_{0} \,\,\, \text{and}\,\,\,E_{\rho}(w^{t},\mathcal{G})<0,\,\,\,\rho\in[\frac{1}{2},1].$$
 We set \(w_{2}:=w^{2t_{2}}\geq 0 \), so that
\begin{equation}\label{eq3.4}
\|w_{2}^{\prime}\|_2^2>2k_{0}\quad\text{and}\quad E_{\rho}(w_{2},\mathcal{G})<0,\quad\forall\rho\in\left[\frac{1}{2},1\right].
\end{equation}
Finally, let \(\Gamma\) and \(c_{\rho}\) be defined as in the statement of the lemma for our choice of \(w_{1}\) and \(w_{2}\); clearly, \(\Gamma\neq \emptyset\), since the path
\[
\gamma_{0}(t):=\frac{\mu^{1/2}}{\|(1-t)w_{1}+tw_{2}\|_2}\left((1-t)w_{1}+tw_{2}\right),\quad  t\in[0,1]
\]
belongs to \(\Gamma\). By continuity and \eqref{eq3.3}--\eqref{eq3.4}, for every \(\gamma \in \Gamma\), there exists a constant \(t_\gamma \in [0,1]\) such that \(\gamma(t_\gamma) \in \partial A_{\mu,k_0}\). Therefore, for any \(\rho \in [\frac{1}{2},1]\), we have, for every \(\gamma \in \Gamma\),
\[
\max_{t \in [0,1]} E_\rho(\gamma(t),\mathcal{G}) \geq E_\rho(\gamma(t_\gamma),\mathcal{G}) \geq \inf_{u \in \partial A_{\mu,k_0}} E_\rho(u,\mathcal{G}) \geq \alpha.
\]
It follows that \(c_\rho \geq \alpha\),  while
\[
\max\{E_\rho(w_1,\mathcal{G}),E_\rho(w_2,\mathcal{G})\} = E_\rho(w_1,\mathcal{G}) < \frac{\alpha}{2},
\]
which completes the proof of the lemma.
\end{proof}

\begin{proof}[Proof of Theorem \ref{th1}]
By Lemma \ref{lem3.1} and Theorem \ref{minimax0}, we deduce that for almost every $\rho\in[\frac{1}{2},1]$, there exists a bounded sequence \(\{u_{\rho,n}\}_{n=1}^{\infty} \subset H^1_\mu(\mathcal{G})\) such that
\begin{equation*}
E_\rho(u_{\rho,n}, \mathcal{G}) \to c_\rho,
\end{equation*}
and
\begin{equation*}
E'_\rho(u_{\rho,n}, \mathcal{G}) + \lambda_{\rho,n} (u_{\rho,n}, \cdot)_2 \to 0 \quad \text{in the dual of } H^1(\mathcal{G}),
\end{equation*}
where
\begin{equation*}
\lambda_{\rho,n} := -\frac{1}{\mu} E'_\rho(u_{\rho,n}, \mathcal{G})[u_{\rho,n}]. 
\end{equation*}
Moreover, there exists a sequence \(\{\zeta_{\rho,n}\} \subset \mathbb{R}^+\) with \(\zeta_{\rho,n}\to 0^+\) such that, if 
\begin{equation*}
 E''_\rho(u_{\rho,n}, \mathcal{G})[\varphi, \varphi] + \lambda_{\rho,n} \|\varphi\|_2^2 < -\zeta_{\rho,n} \|\varphi\|^2 
\end{equation*}
holds for any \( \varphi \in W_n \setminus \{0\} \) in a subspace \( W_n \) of \( T_{u_{\rho,n}} H^1_\mu(\mathcal{G}) \), then   \(\dim W_n \leq  1 \). In addition, as noted in \cite[Remark 1.8]{Bort0}, since $u\in H^{1}_{\mu}(\mathcal{G})\implies|u|\in H^{1}_{\mu}(\mathcal{G})$, $w_{1},w_{2}\geq 0$ (see Lemma \ref{lem3.1}), the map $u\mapsto|u|$ is continuous, and $E_{\rho}(u)=E_{\rho}(|u|)$, it is possible to choose $\{u_{\rho,n}\}$ with the property that $u_{\rho,n}\geq 0$ on $\mathcal{G}$. Moreover, the monotonicity of $c_\rho$ as a function of $\rho\in [\frac{1}{2},1]$ implies that $\{c_\rho\}$ is bounded, with $c_1\leq c_\rho \leq c_\frac{1}{2}$.

By Lemma \ref{lamuna+}, up to a subsequence, $\displaystyle \lambda_{\rho}:= \lim_{n \to \infty}\lambda_{\rho,n}\geq 0$. Moreover, by  Lemmas \ref{lemurhosolve} and \ref{lemurhoconverge}, there exists a solution $u_\rho\in H^1(\G)$ to \eqref{eqpdnomu} such that $u_{\rho,n}(x) \to u_\rho(x)$ for all $x \in \G$,  and
\begin{equation*}
\int_{\mathcal{G}}|(u_{\rho,n}-u_{\rho})^{\prime}|^{2}\,dx+\lambda_{\rho}\int_{\mathcal{G}}|u_{\rho,n}-u_{\rho}|^{2}\,dx \to 0, \quad \text{as }n \to \infty.
\end{equation*} 
Since $u_{\rho,n}\geq 0$ on $\mathcal{G}$, we have $u_\rho \geq 0$ on $\mathcal{G}$.

Suppose, by contradiction, that $u_{\rho}\equiv 0$. Then
\begin{equation}\label{3.17}
\int_{\mathcal{G}}|u_{\rho,n}^{\prime}|^{2}\,dx+\lambda_{\rho}\int_{\mathcal{G}}|u_{\rho,n}|^{2}\,dx\to 0,\quad \text{as }n \to \infty.
\end{equation}
If $\lambda_{\rho}>0$, this is impossible since $\|u_{\rho,n}\|^{2}_2=\mu>0$. If $\lambda_{\rho}=0$, then, by the Gagliardo-Nirenberg inequality \eqref{eqgn}, \eqref{3.17} contradicts the fact that  $c_{\rho}>0$. Hence, $u_{\rho}\not\equiv 0$.
By the vertex conditions and the uniqueness theorem for ODEs, we conclude that $u_{\rho}>0$ in $\mathcal{G}$. Indeed, suppose that there exists $x_0\in\G$ such that $u_\rho(x_0)=0$. If $x_0$ lies in the interior of an edge, then the strong
maximum principle implies that $u_\rho$ vanishes identically on that edge.
In particular, $u_\rho$ vanishes at an endpoint of that edge. Therefore, it
is enough to consider the case $x_0\in\V$. Since $u_\rho\geq0$ on $\G$ and $u_\rho(x_0)=0$, one has $u_{\rho,\e}'(x_0)\geq0$ for each $ \e\succ x_0$
where $(u_{\rho,\e})'(x_0)$ stands for $u_{\rho}|_\e'(0)$ or $-u_{\rho}|_\e'(\ell_\e)$, according to whether the vertex $x_0$ is identified with $0$ or $\ell_\e$. Moreover, by \eqref{eqpdnomu}, we obtain
 $$\sum_{\e\succ x_0}u_{\rho,\e}'(x_0)=0,
 $$
 and thus, $u_{\rho,\e}'(x_0)=0$ for all $\e \succ x_0$. Hence, by the uniqueness of solution to second order ODEs (see, e.g., {\cite[Theorem 10 in Section 3.1]{AD}}), $u_\rho \equiv 0$ on any edge containing $x_0$, and repeating this argument over finitely many edges (since $\mathcal{G}$ has finitely many vertices and edges), we deduce that $u_{\rho}\equiv 0$ on $\mathcal{G}$, which leads to a contradiction.


Next, we claim that $\lambda_\rho > 0$. If $\lambda_\rho = 0$, consider an unbounded edge $\e_1 \in \E$ identified with $[0, +\infty)$. Then,  $u_\rho$ is $C^2$ on $[0, +\infty)$ and satisfies $u_\rho''(x) = 0$ for all $x \in (0, +\infty)$. It follows from \cite[Remark 4 in Section 3.3]{AD} that $u_\rho > 0$ on $[0, +\infty)$ and
\[
u_\rho(x) = u_\rho(0) + u_\rho'(0)x,
\]
which leads to a contradiction with $u_\rho \in L^2(\mathcal{G})$. Hence the claim holds. 

By Lemma \ref{lemurhoconverge}, we conclude that $(u_\rho,\lambda_\rho) \in H^1(\G)\times (0,\infty)$ is a solution to equation \eqref{eqpdrho} and satisfies $E_{\rho}(u_{\rho},\G)=c_{\rho}$.

Finally, there exists a subsequence $\{\rho_n\} \subset [\frac{1}{2},1]$ with $\rho_n \to 1^-$ such that each $(u_{\rho_n},\lambda_{\rho_n}) \in H^1(\G) \times (0, +\infty)$ solves 
\begin{align*}
\begin{cases} 
u_{\rho_n}'' = \lambda_{\rho_n} u_{\rho_n} & \text{on every }\e \in \E, \\ 
\int_\G\abs{u_{\rho_n} }^2\, dx = \mu & \\
\displaystyle\sum_{\e \succ \vv} u_{\rho_n,\e}' (\vv)= -\rho_n|u_{\rho_n}(\vv)|^{p-2}u_{\rho_n}(\vv)  & \text{at every } \vv \in \V_0,\\
\displaystyle\sum_{\e \succ \vv} u_{\rho_n,\e}' (\vv)= 0  & \text{at every } \vv \in \V \setminus \V_0
\end{cases}
\end{align*}
with $E_{\rho_n}(u_{\rho_n},\G)=c_{\rho_n}$, where $u_{\rho_n,\e}'(\vv)$ stands for $u_{\rho_n}|_\e'(0)$ or $-u_{\rho_n}|_\e'(\ell_\e)$, according to whether the vertex $\vv$ is identified with $0$ or $\ell_\e$. In view of $c_1\leq c_{\rho_n}\leq c_{1/2}$ and Lemma \ref{lembu}, $\{u_{\rho_{n}}\} \subset H^{1}(\G)$ is a bounded sequence satisfying the assumptions \eqref{PS1} and \eqref{PS2}. 
By repeating the previous arguments, we conclude that, up to a subsequence, $\displaystyle\lambda =\lim_{n \to \infty}\lambda_{\rho_{n}}>0$, $\{u_{\rho_{n}}\} $ converges to some $u \in H_\mu^1(\mathcal{G})$ with $u>0$ on $\G$ and $(u, \lambda) \in H_\mu^1(\mathcal{G}) \times (0,+\infty)$ is a solution to \eqref{eqpd}.
\end{proof}



\section{Proof of Theorem \ref{th2}}\label{sect5}
Throughout this section, we consider the metric graph $\G_M=(\V,\E_M)$ defined in the Introduction and the nonempty nonlinear defect set $\V_0=\{\vv_1,\ldots,\vv_q\}$. We prove the multiplicity result for solutions of problem \eqref{eqpdm}. 

 For each $i=1,\ldots,q$, we denote the degree of the vertex $\vv_i$ in
$\G_M$ by
$$
d_M(\vv_i):=\#\{\e\in\E_M:\e\succ\vv_i\}.
$$
For $N=2,\ldots,q+1$, let $K_{M,N-1}:=\max_{1\leq i \leq N-1} d_M(\vv_i)$. Since $\G$ has finitely many edges, we may
choose $\ell>0$ sufficiently small independent of $M$ such that the $q$ sets
$$
\bigcup_{\e\succ\vv_i, e \in E_M}
\left(
\e\cap\left[0,\frac{\ell}{d_M(\vv_i)}\right]
\right),
\qquad i=1,\ldots,q,
$$ 
are pairwise disjoint and contain no vertex of $\G_M$ except $\vv_i$, where
$\vv_i$ is identified with $0$ on every incident edge.


For $N=2,\ldots,q+1$, set $k=N-1$. We shall construct maps
$$
\gamma_{0,N},\gamma_{1,N}:{\mathbb S}^{N-2}\to H^1_\mu(\G_M)
$$
and then use the linking class appearing in Theorem \ref{minimax}, namely
\begin{equation*}
\Gamma_N:=\left\{\gamma\in C([0,1]\times{\mathbb S}^{N-2},H^1_\mu(\G_M)):\begin{array}{l}
\gamma(t,\cdot)\text{ is odd for every }t\in[0,1],\\
\gamma(0,\cdot)=\gamma_{0,N},\quad \gamma(1,\cdot)=\gamma_{1,N}
\end{array}\right\}.
\end{equation*}
The corresponding min-max level is
\begin{equation*}
c_{\rho,\mu}^N:=\inf_{\gamma\in\Gamma_N}\max_{(t,a)\in[0,1]\times{\mathbb S}^{N-2}}E_\rho(\gamma(t,a),\G_M).
\end{equation*}
In what follows, we present two lemmas that construct the orthogonal families required in Theorem \ref{minimax}.
\begin{lemma}\label{lem lower endpoint}
For any $\mu>0$, any $M\geq 2$ and any $N=2,\ldots,q+1$, there exist functions \(\{\varphi_{1},\varphi_{2},\ldots,\varphi_{N-1}\}\) such that the following properties hold for all \(1\leq i,j\leq N-1\):
	\begin{enumerate}[label=(\roman*)]
		\item\(\varphi_{i}\in H^1_\mu(\G_M)\), \(\|\varphi^{\prime}_{i}\|_2=1\), and \(E_{\rho}(\varphi_{i},\mathcal{G}_M)=1/2, \quad \forall \rho \in [1/2,1]\);\label{lem4.31}
		\item $\operatorname{supp}(\varphi_{i}) \subset \{x \in \mathcal{H}: \operatorname{dist}(x,\vv_\mathcal{H}) \geq 2\ell\}$ for some half-line $\mathcal{H} \in \E_M$ and the endpoint $\vv_\mathcal{H}$ of $\mathcal{H}$, and \(\operatorname{supp}(\varphi_{i})\cap\operatorname{supp}(\varphi_{j})=\varnothing\) whenever \(i\neq j\);\label{lem4.32}
		\item if \(a=(a_1,...,a_{N-1})\in\mathbb{S}^{N-2}\),  then\\ 
        $$\big{\|}\big{(}\sum_{i=1}^{N-1}a_{i}\varphi_{i}\big{)}^{\prime}\big{\|}_2=1 \text{ and } E_{\rho}\big{(}\sum_{i=1}^{N-1}a_{i}\varphi_{i},\,\mathcal{G}_M\big{)}=1/2, \quad \forall \rho \in [1/2,1].$$\label{lem4.33}
\end{enumerate}
\end{lemma}

\begin{proof}
Let \(\varphi\in C_c^\infty(\mathbb{R})\) satisfy \(\operatorname{supp}\varphi\subset(0,1)\) and \(\|\varphi\|_{L^2(\mathbb{R})}^2=\mu\). 
	Viewing \(\varphi\) as a function in \(H^{1}(\mathcal{G}_M)\) whose support is contained in a half-line $\mathcal{H}$, identified with \([0,\infty)\), we use \eqref{scal} to define
	\[
	\varphi_{1}(x):=\tau^{1/2}\varphi\left(\tau (x-2\ell)\right)\quad\text{with}\quad\tau:=\frac{1}{\|\varphi^{\prime}\|_2}.
	\]
	The function \(\varphi_{1}\) satisfies \ref{lem4.31}. Indeed, for any \(t>0\), \(\|\varphi^{t}\|_2^{2}=\|\varphi\|_2^{2}=\mu\), and a direct calculation yields
	\[
	\|\varphi^{\prime}_{1}\|^{2}_2=\tau^{2}\|\varphi^{\prime}\|^{2}_2=1.
	\]
	Finally, since \(\varphi_{1}\) is supported in the half-line, we have
	\[
	E_{\rho}(\varphi_{1},\mathcal{G}_M)=\tfrac{1}{2}\|\varphi^{\prime}_{1}\|^{2}_2=\frac{1}{2}.
	\]
	Now define, for \( 2\leq i\leq N-1\),
	\[
	\varphi_{i}(x):=\varphi_{1}\left(x-\frac{i-1}{\tau}\right).
	\]
	Since \(\varphi_{i}\) are translations of \(\varphi_{1}\), they still satisfy \ref{lem4.31}. Also, observe that, by definition, \(\operatorname{supp}(\varphi_{i})\subset\big{(}2\ell+\frac{i-1}{\tau},2\ell+\frac{i}{\tau}\big{)}\), and so they all have disjoint compact supports. This proves \ref{lem4.32}.

 Finally, observe that for any \(a\in\mathbb{S}^{N-2}\)
	\[
	\left\|\left(\sum_{i=1}^{N-1}a_{i}\varphi_{i}\right)^{\prime}\right\|^{2}_2=\sum_{i=1}^{N-1}a_{i}^{2}\,\|\varphi^{\prime}_{i}\|^{2}_2=1,
	\]
	from which \ref{lem4.33} follows. This proves the lemma.
\end{proof}
For any $N=2,\ldots,q+1$ and any $M \geq 2$, set
$$
\beta_{N,\mu}(M):=2^\frac{p}{4-p}M^\frac{p(N-1)}{p-4}\mu^{p/(8-2p)}.
$$
\begin{lemma}\label{lem controlled upper endpoint}
There exists $L_0=L_0(p,q,\ell)>0$, independent of $M,N,\rho$ and $\mu$, such that, for any $0<\mu<L_0^\frac{p-4}{p-2}$, any $M\geq 2$ and any $N=2,\ldots,q+1$, 
 there exist functions \(\{\bar{\varphi}_{1},\bar{\varphi}_{2},\ldots,\bar{\varphi}_{N-1}\}\) such that for any \(1\leq i,j\leq N-1\), we have
\begin{enumerate}[label=(\roman*)]
		\item\(\bar{\varphi}_{i}\in H^1_\mu(\G_M)\), \( \|\overline{\varphi}^{\prime}_{i}\|_2 =\frac{\sqrt{3}L_{0}}{\ell}
K^{\frac{p}{p-4}}_{M,N-1}
\mu^\frac{p}{8-2p} \geq 2\beta_{N,\mu}(M) \);\label{lem7.51}
		\item $\operatorname{supp}(\bar{\varphi}_{i}) \subset \{x\in \G_M:\operatorname{dist}(x,\vv_i)\leq\ell\}$ and \(\operatorname{supp}(\bar{\varphi}_{i})\cap\operatorname{supp}(\bar{\varphi}_{j})=\varnothing\) whenever \(i\neq j\);\label{lem7.52}
		\item  if \(a=(a_1,...,a_{N-1})\in\mathbb{S}^{N-2}\),  then $$\bigl{\|}\bigl{(}\sum_{i=1}^{N-1}a_{i}\overline{\varphi}_{i}\bigr{)}^{\prime}\bigr{\|}_2=\frac{\sqrt{3}L_0}{\ell}
K^{\frac{p}{p-4}}_{M,N-1}
\mu^\frac{p}{8-2p} \geq 2\beta_{N,\mu}(M) \quad  \text{and}\quad E_{\rho}\bigl{(}\sum_{i=1}^{N-1}a_{i}\overline{\varphi}_{i},\,\mathcal{G}_{M}\bigr{)} \leq 0, \quad \forall \rho \in[1/2,1].$$\label{lem7.53}
\end{enumerate}
\end{lemma}
\begin{proof}
Define $u_{s_i}\in H^1(\G_M)$ by
$$
u_{s_i}(x)=
\begin{cases}
s_i\left(\dfrac{\ell}{d_M(\vv_i)}-x\right),
&\text{on } \e\cap\left[0,\dfrac{\ell}{d_M(\vv_i)}\right],\ \e\succ\vv_i,\\
0,
&\text{elsewhere},
\end{cases}
$$
where $\vv_i$ is identified with $0$ on each edge $\e\succ\vv_i$.

We choose $s_i>0$ such that
$$
\|u_{s_i}\|_2^2=\mu,
$$
and set
$$
\phi_i:=u_{s_i}.
$$
A direct computation gives
$$
\|\phi_i'\|_2^2
=
\frac{3\mu \left(d_M(\vv_i)\right)^2}{\ell^2},
$$
and
$$
\phi_i(\vv_i)
=
s_i\frac{\ell}{d_M(\vv_i)}
=
\sqrt{\frac{3\mu}{\ell}}.
$$
Let $L>0$ be a constant to be chosen later. For $i=1,\ldots,N-1$, define
$$
t_i
:=
L\frac{K^{\frac{p}{p-4}}_{M,N-1}}{d_M(\vv_i)}
\mu^{-\frac{p-2}{p-4}}.
$$
For $0<\mu<L^\frac{p-4}{p-2}$, we have
$t_i\geq1$ for all $i=1,\ldots,N-1$.
Define
$$
\bar{\varphi}_i:=\phi_i^{t_i},
\qquad i=1,\ldots,N-1,
$$
where
$$
\phi_i^{t_i}(x):=t_i^{1/2}\phi_i(t_ix)
$$
on every edge incident at $\vv_i$, with $\vv_i$ identified with $0$.

Then
$$
\bar{\varphi}_i\in H^1_\mu(\G_M),
\qquad i=1,\ldots,N-1.
$$
Moreover,
$$
\operatorname{supp}(\bar{\varphi}_i)
\subset
\bigcup_{\e\succ\vv_i}
\left(
\e\cap\left[0,\frac{\ell}{t_id_M(\vv_i)}\right]
\right)
\subset
\bigcup_{\e\succ\vv_i}
\left(
\e\cap\left[0,\frac{\ell}{d_M(\vv_i)}\right]
\right),
$$
and hence the supports of
$\bar{\varphi}_1,\ldots,\bar{\varphi}_{N-1}$ are pairwise disjoint. This proves \ref{lem7.52}.
For $a=(a_1,\ldots,a_{N-1})\in\mathbb S^{N-2}$, set
$$
u_a:=\sum_{i=1}^{N-1}a_i\bar{\varphi}_i.
$$
Then $u_a\in H^1_\mu(\G_M)$.

We compute
$$
\|\bar{\varphi}_i'\|_2^2
=
t_i^2\frac{3\mu \left(d_M(\vv_i)\right)^2}{\ell^2}= \frac{3L^2}{\ell^2}
K^{\frac{2p}{p-4}}_{M,N-1}
\mu^\frac{p}{4-p}.
$$
This proves the equality in \ref{lem7.51}; the lower bound by $2\beta_{N,\mu}(M)$ follows from the choice of $L_0$ below.
Consequently,
$$
\|u_a'\|_2^2
=\sum_{i=1}^{N-1}a_i^2\|\bar{\varphi}_i'\|_2^2= \frac{3L^2}{\ell^2}
K^{\frac{2p}{p-4}}_{M,N-1}
\mu^\frac{p}{4-p}.
$$
Moreover,
$$
|\bar{\varphi}_i(\vv_i)|^p
=
t_i^{p/2}
\left(\frac{3\mu}{\ell}\right)^{p/2}=
3^{p/2}\ell^{-p/2}L^{p/2}
\left(\frac{K_{M,N-1}}{d_M(\vv_i)}\right)^{p/2}
K^{\frac{2p}{p-4}}_{M,N-1}\mu^{\frac{p}{4-p}}.
$$
It follows from $K_{M,N-1} \geq d_M(\vv_i)$ that
$$
\sum_{\vv\in\V_0}|u_a(\vv)|^p
\geq
3^{p/2}\ell^{-p/2}L^{p/2}
K^{\frac{2p}{p-4}}_{M,N-1}\mu^{\frac{p}{4-p}}
\sum_{i=1}^{N-1}|a_i|^p.
$$
Since $a\in\mathbb S^{N-2}$ and $N-1\leq q$, we have
$$
\sum_{i=1}^{N-1}|a_i|^p
\geq
(N-1)^{1-\frac p2}
\geq
q^{1-\frac p2}.
$$
Hence,
$$
\sum_{\vv\in\V_0}|u_a(\vv)|^p
\geq
3^{p/2}q^{1-\frac p2}\ell^{-p/2}L^{p/2}
K^{\frac{2p}{p-4}}_{M,N-1}\mu^{\frac{p}{4-p}}.
$$
Using $\rho\geq\frac12$, we obtain
$$
\begin{aligned}
E_\rho(u_a,\G_M)
&=
\frac12\|(u_a)'\|_2^2
-
\frac{\rho}{p}
\sum_{\vv\in\V_0}|u_a(\vv)|^p
\\
&\leq
\left(
\frac{3}{2\ell^2}L^2
-
\frac{3^{p/2}q^{1-\frac p2}}{2p\ell^{p/2}}L^{p/2}
\right)
K^{\frac{2p}{p-4}}_{M,N-1}\mu^{\frac{p}{4-p}}.
\end{aligned}
$$
Since $p>4$, we may choose $L=L_0(p,q,\ell)>0$ sufficiently large so that
\[
\frac{3}{\ell^2}L_0^2
-
\frac{3^{p/2}q^{1-\frac p2}}{p\ell^{p/2}}L_0^{p/2}
\leq0
\]
and
\[
\frac{\sqrt{3}L_0}{\ell}\,2^{\frac{p}{p-4}}\geq2.
\]
The first condition ensures the required nonpositivity of the endpoint energy. Moreover, since $K_{M,N-1}\geq M^{N-1}$, the second condition gives, for every $i=1,\ldots,N-1$,
\[
\|\bar{\varphi}_i'\|_2
=\frac{\sqrt{3}L_0}{\ell}K_{M,N-1}^{\frac{p}{p-4}}
\mu^{\frac{p}{8-2p}}
\geq2\beta_{N,\mu}(M).
\]
Since the functions have pairwise disjoint supports and all their derivative norms are equal, the same computation yields
\[
\|u_a'\|_2
=\frac{\sqrt{3}L_0}{\ell}K_{M,N-1}^{\frac{p}{p-4}}
\mu^{\frac{p}{8-2p}}
\geq2\beta_{N,\mu}(M).
\]
With this choice
$$
E_\rho(u_a,\G_M)\leq 0,
$$ 
from which \ref{lem7.53} follows.
This completes the proof of Lemma \ref{lem controlled upper endpoint}.
\end{proof}
Define
$$
\gamma_{0,N}:\mathbb S^{N-2}\to H^1_\mu(\G_M)\quad\text{by}\quad\gamma_{0,N}(a_{1}, \ldots,a_{N-1})=\sum_{i=1}^{N-1}a_{i}\varphi_{i}
$$
where $\{\varphi_{1},\varphi_{2},\ldots,\varphi_{N-1}\}$ is given by Lemma \ref{lem lower endpoint}, and denote
$$
\gamma_{1,N}:\mathbb S^{N-2}\to H^1_\mu(\G_M)\quad\text{by}\quad \gamma_{1,N} (a_{1}, \ldots,a_{N-1})=\sum_{i=1}^{N-1}a_{i}\bar{\varphi}_{i}
$$
where $\{\bar{\varphi}_{1},\bar{\varphi}_{2},\ldots,\bar{\varphi}_{N-1}\}$ is defined in Lemma \ref{lem controlled upper endpoint}. Define
\begin{equation*}
\Gamma_N:=\left\{\gamma\in C([0,1]\times{\mathbb S}^{N-2},H^1_\mu(\G_M)):\begin{array}{l}
\gamma(t,\cdot)\text{ is odd for every }t\in[0,1],\\
\gamma(0,\cdot)=\gamma_{0,N},\quad \gamma(1,\cdot)=\gamma_{1,N}
\end{array}\right\}.
\end{equation*}
By \cite[Lemma 2.8]{Ca}, we know that $\Gamma_N$ is nonempty. 
The corresponding min-max level is
\begin{equation*}
c_{\rho,\mu}^N:=\inf_{\gamma\in\Gamma_N}\max_{(t,a)\in[0,1]\times{\mathbb S}^{N-2}}E_\rho(\gamma(t,a),\G_M).
\end{equation*}
We first derive an upper bound for $c_{\rho,\mu}^N$.
\begin{lemma}\label{lem upper bound cN}
For any $0<\mu<\min\left\{K_{M,N-1}^2
\left(\frac{3L_0^2}{\ell^2}\right)^{\frac{p-4}{p}},L_0^\frac{p-4}{p-2}\right\}$, any $M\geq 2$, and any
$N=2,\ldots,q+1$, one has
$$
c_{\rho,\mu}^N
\leq
\frac{3L_0^2}{2\ell^2}K^{\frac{2p}{p-4}}_{M,N-1}
\mu^\frac{p}{4-p}, \quad \rho\in [1/2,1].
$$
\end{lemma}

\begin{proof}
By Lemmas \ref{lem lower endpoint} and \ref{lem controlled upper endpoint}, if $0<\mu<L_0^{\frac{p-4}{p-2}}$,
$$
\|\gamma_{0,N}(a)'\|_2=1,
\qquad
E_\rho(\gamma_{0,N}(a),\G_M)=\frac12,
$$
and
$$
\|\gamma_{1,N}(a)'\|_2^2
\leq
\frac{3L_0^2}{\ell^2}
K^{\frac{2p}{p-4}}_{M,N-1}
\mu^\frac{p}{4-p}
$$
and
$$
E_\rho(\gamma_{1,N}(a),\G_M)\leq 0.
$$
Since the supports of $\{{\varphi}_{1},{\varphi}_{2},\ldots,{\varphi}_{N-1}\}$ are disjoint from the supports of $\{\bar{\varphi}_{1},\bar{\varphi}_{2},\ldots,\bar{\varphi}_{N-1}\}$,
$$
(\gamma_{0,N}(a),\gamma_{1,N}(a))_{L^2(\G_M)}=0,
\qquad 
\forall a\in\mathbb S^{N-2}.
$$
Define
$$
\gamma_N(s,a)
:=
\frac{(1-s)\gamma_{0,N}(a)+s\gamma_{1,N}(a)}
{\sqrt{(1-s)^2+s^2}},
\qquad
(s,a)\in[0,1]\times\mathbb S^{N-2}.
$$
Since
$$
\|\gamma_{0,N}(a)\|_2^2
=
\|\gamma_{1,N}(a)\|_2^2
=
\mu
$$
and the two functions are $L^2$-orthogonal, we have
$$
\gamma_N(s,a)\in H^1_\mu(\G_M).
$$
Moreover, $\gamma_N(s,\cdot)$ is odd, and
$$
\gamma_N(0,a)=\gamma_{0,N}(a),
\qquad
\gamma_N(1,a)=\gamma_{1,N}(a).
$$
Therefore
$$
\gamma_N\in\Gamma_N.
$$
By disjointness of the supports, we have
$$
\begin{aligned}
\|\gamma_N(s,a)'\|_2^2
&=
\frac{
(1-s)^2\|\gamma_{0,N}(a)'\|_2^2
+
s^2\|\gamma_{1,N}(a)'\|_2^2
}
{(1-s)^2+s^2}
\\
&\leq
\max\left\{
1,
\frac{3L_0^2}{\ell^2}K^{\frac{2p}{p-4}}_{M,N-1}
\mu^\frac{p}{4-p}
\right\}.
\end{aligned}
$$
For $0<\mu<\min\left\{K^{2}_{M,N-1}\left(\frac{3L_0^2}{\ell^2}\right)^\frac{p-4}{p},L_0^\frac{p-4}{p-2}\right\}$, we get
$$
E_\rho(\gamma_N(s,a),\G_M)
\leq
\frac12\|\gamma_N(s,a)'\|_2^2
\leq
\frac{3L_0^2}{2\ell^2}K^{\frac{2p}{p-4}}_{M,N-1}
\mu^\frac{p}{4-p}.
$$
Taking the maximum over $(s,a)\in[0,1]\times\mathbb S^{N-2}$ and then the
infimum over $\Gamma_N$, we obtain
$$
c_{\rho,\mu}^N
\leq
\frac{3L_0^2}{2\ell^2}K^{\frac{2p}{p-4}}_{M,N-1}
\mu^\frac{p}{4-p},
$$
which completes the proof of Lemma \ref{lem upper bound cN}.
\end{proof}
We now derive a lower bound for $c^N_{\rho,\mu}$.
\begin{lemma}\label{lem-trace-K}
If $M \geq 2$, then for every $N=2,\ldots,q+1$, and every
\(u\in H^1_\mu(\G_M)\),
\begin{equation}\label{trace-tail}
\sum_{i=N-1}^{q}|u(\vv_i)|^p\leq 2^{p/2+1}M^{(p-pN)/2}\mu^{p/4}\norm{u'}_2^{p/2}.
\end{equation}


%where $C_{\rm tr}:=2^{p/2+1}$ is independent of $M,N,\rho$ and $\mu$.
\end{lemma}
\begin{proof}
Let $h_{i,1},\ldots,h_{i,M^{i}}$ be the new half-lines attached at $\vv_i$. By the Gagliardo--Nirenberg inequality \eqref{eqgn} (or the one-dimensional trace inequality on the half-line), we have
$$
|u(\vv_i)|^2\leq2\norm{u}_{L^2(h_{i,j})}\norm{u'}_{L^2(h_{i,j})}.
$$
Summing this inequality over $j=1,\ldots,M^{i}$ and using the Cauchy--Schwarz inequality, we obtain
\begin{equation}\label{traceKi}
M^{i}|u(\vv_i)|^2\leq2\left(\sum_{j=1}^{M^{i}}\norm{u}_{L^2(h_{i,j})}^2\right)^{1/2}\left(\sum_{j=1}^{M^{i}}\norm{u'}_{L^2(h_{i,j})}^2\right)^{1/2}\leq2\mu^{1/2}\norm{u'}_2.
\end{equation}
For $i\geq N-1$, we have
$$
M^{-ip/2}=M^{(p-pN)/2}M^{-(i-N+1)p/2}.
$$
Thus, for $M\geq2$,
$$
\sum_{i=N-1}^{q}M^{-ip/2}\leq M^{(p-pN)/2}\sum_{i=0}^{q+1-N}M^{-ip/2}\leq2M^{(p-pN)/2}.
$$
Combining this with \eqref{traceKi} gives \eqref{trace-tail}.
\end{proof}
For $N=2,\ldots,q+1$, we define the closed subspace
\begin{equation}\label{zero-trace-W}
Z_{N-2}^M
:=
\left\{
u\in H^1(\G_M):
u(\vv_1)=\cdots=u(\vv_{N-2})=0
\right\},
\end{equation}
with $Z_0^M=H^1(\G_M)$, and define $W_{N-2}^M:=\left(Z_{N-2}^M\right)^\perp$.
Then $\dim W_{N-2}^M\leq N-2$ and $\left(W_{N-2}^M\right)^\perp=Z_{N-2}^M$. Indeed,  since the maps
$u\mapsto u(\vv_i)$, $i=1,\ldots,N-2$,
are continuous linear functionals on $H^1(\G_M)$, by the
Riesz representation theorem, there exist functions
$R_i^M\in H^1(\G_M)$ such that
\[
u(\vv_i)=\left(u,R_i^M\right)_{H^1(\G_M)},
\qquad
\forall u\in H^1(\G_M).
\]
Hence
\[
Z_{N-2}^M
=
\left(\operatorname{span}\{R_1^M,\ldots,R_{N-2}^M\}\right)^\perp.
\]
Consequently,
\[
\dim \left(W_{N-2}^M\right) =\dim\left(Z_{N-2}^M\right)^\perp\leq N-2.
\]

Define
\begin{equation}\label{BNmu}
B_{N,\mu}^M:=\{u\in (W_{N-2}^M)^\perp\cap H^1_\mu(\G_M):\norm{u'}_2=\beta_{N,\mu}(M)\}.
\end{equation}



\begin{lemma}\label{lem-window-lower}
For any $ \mu >0$, any $M \geq 2$, and any $N=2,\ldots,q+1$,
\begin{equation}\label{lower-shell}
\inf_{u\in B_{N,\mu}^M}E_\rho(u,\G_M)\geq \frac{(p-4)2^{\frac{3p-4}{4-p}}}{p}M^{\frac{2p(N-1)}{p-4}}\mu^{\frac{p}{4-p}},\quad \forall \rho\in[1/2,1].
\end{equation}
Moreover, for any $0<\mu<\min\left\{\frac{M^{2N-2}}{4}\left( \frac{p-4}{p} \right)^{\frac{p-4}{p}},L_0^\frac{p-4}{p-2}\right\}$, any $M\geq 2$, and any
$N=2,\ldots,q+1$, one has
\begin{equation}\label{eqlem55}
\begin{aligned}
c_{\rho,\mu}^N
\ge\inf_{u\in B_{N,\mu}^M}E_\rho(u,\G_M) &\geq  \frac{(p-4)2^{\frac{3p-4}{4-p}}}{p}M^{\frac{2p(N-1)}{p-4}}\mu^{\frac{p}{4-p}}\\
&>\frac{1}{2}= \max_{s \in \mathbb{S}^{N-2}}\max\{E_\rho(\gamma_{0,N}(s),\G_M),E_\rho(\gamma_{1,N}(s),\G_M)\},\quad \forall \rho\in[1/2,1].
\end{aligned}
\end{equation}
\end{lemma}

\begin{proof}
Let $u\in (W_{N-2}^M)^\perp\cap H^1_\mu(\G_M)$. By \eqref{zero-trace-W},  $u(\vv_1)=\cdots=u(\vv_{N-2})=0$. Hence, using \eqref{trace-tail},
$$
\sum_{\vv\in\V_0}|u(\vv)|^p=\sum_{i=N-1}^{q}|u(\vv_i)|^p\leq 2^{p/2+1}M^{(p-pN)/2}\mu^{p/4}\norm{u'}_2^{p/2}.
$$
Since $\rho\leq1$, we get
$$
E_\rho(u,\G_M)\geq\frac12\norm{u'}_2^2-\frac{2^{p/2+1}M^{(p-pN)/2}}{p}\mu^{p/4}\norm{u'}_2^{p/2}.
$$
On $B_{N,\mu}^M$, we have
$$
\begin{aligned}
    E_\rho(u,\G_M)&\geq\mu^{p/(4-p)}\left(\frac12\left(\frac{1}{2^{p/2}M^{(p-pN)/2}}\right)^{\frac{4}{p-4}}-\frac{2^{p/2+1}M^{(p-pN)/2}}{p}\left(\frac{1}{2^{p/2}M^{(p-pN)/2}}\right)^{\frac{p}{p-4}}\right)\\
&\geq \frac{p-4}{2p}\mu^{p/(4-p)}\left(\frac{1}{2^{p/2}M^{(p-pN)/2}}\right)^{\frac{4}{p-4}}\\
& \geq \frac{(p-4)2^{\frac{3p-4}{4-p}}}{p} M^{\frac{2p(N-1)}{p-4}}\mu^{\frac{p}{4-p}}.
\end{aligned}
$$

We now apply \cite[Theorem 2.10]{Ca} with the choices \(\Phi=E_{\rho}(\cdot,\mathcal{G}_M), d=N-2, J(u)=\|u^{\prime}\|_2, \beta=\beta_{N,\mu}(M)\), and \(W=W^M_{N-2}\). Indeed,
for any $$0<\mu<\min\left\{\frac{M^{2N-2}}{4}\left( \frac{p-4}{p} \right)^{\frac{p-4}{p}},L_0^\frac{p-4}{p-2}\right\},$$ 
since $\frac{M^{2N-2}}{4}\left( \frac{p-4}{p} \right)^{\frac{p-4}{p}}<\frac{M^{2N-2}}{4}$,  we have 
\[
\|\gamma_{0,N}(a)'\|_2=1<2^\frac{p}{4-p}M^\frac{p(N-1)}{p-4}\mu^{p/(8-2p)}= \beta_{N,\mu}(M), \quad \|\gamma_{1,N}(a)'\|_2\geq2\beta_{N,\mu}(M),
\]
and
$$
\begin{aligned}
\inf_{u\in B_{N,\mu}^M}E_\rho(u,\G_M)&\geq  \frac{(p-4)2^{\frac{3p-4}{4-p}}}{p} M^{\frac{2p(N-1)}{p-4}}\mu^{\frac{p}{4-p}}\\
&> \frac{1}{2}= \max_{s \in \mathbb{S}^{N-2}}\max\{E_\rho(\gamma_{0,N}(s),\G_M),E_\rho(\gamma_{1,N}(s),\G_M)\},
\end{aligned}
$$
and the space $W_{N-2}^M$ has dimension at most $N-2$. Thus, the endpoint conditions $J(\gamma_{0,N}(a))<\beta_{N,\mu}(M)<J(\gamma_{1,N}(a))$, the strict energy separation from the shell $B_{N,\mu}^M$, and the dimension condition $\dim W_{N-2}^M\leq N-2$ required in \cite[Theorem 2.10]{Ca} are all satisfied. Consequently, $c_{\rho,\mu}^N\geq\inf_{u\in B_{N,\mu}^M}E_\rho(u,\G_M)$, and \eqref{eqlem55} follows. This completes the proof of the lemma.
\end{proof}

Set
$$
\begin{aligned}
    \mu_0:=\min{\left\{4\left(\frac{3L_0^2}{\ell^2}\right)^\frac{p-4}{p},\left( \frac{p-4}{p} \right)^{\frac{p-4}{p}},L_0^\frac{p-4}{p-2}\right\}}.
\end{aligned}
$$
By construction, $\mu_0$ depends only on $p,q$ and $\ell$. For $M \geq2$ and $N=2,\ldots,q+1$,
$$
    \mu_0 \leq \min\left\{K^{2}_{M,N-1}\left(\frac{3L_0^2}{\ell^2}\right)^\frac{p-4}{p},\frac{ M^{2N-2}}{4} \left( \frac{p-4}{p} \right)^{\frac{p-4}{p}},L_0^\frac{p-4}{p-2}\right\}.
$$
\begin{lemma}\label{prop-windows}
There exists an integer \(M_0=M_0(p,q,\ell,\G,\V_0)\geq2\), independent of \(\mu\), such that,
for every \(M\geq M_0\) and every \(0<\mu<\mu_0\),
\begin{equation}\label{sep-levels}
c_{1,\mu}^2\leq c_{1/2,\mu}^2
<
c_{1,\mu}^3\leq c_{1/2,\mu}^3
<
\cdots
<
c_{1,\mu}^{q+1}\leq c_{1/2,\mu}^{q+1}.
\end{equation}
\end{lemma}
\begin{proof}
Since $0< \mu < \mu_0$, by Lemmas \ref{lem upper bound cN} and \ref{lem-window-lower}, for $N=2,\ldots,q+1$,
\begin{equation}\label{ep-levels}
 \frac{(p-4)2^{\frac{3p-4}{4-p}}}{p}M^{\frac{2p(N-1)}{p-4}}\mu^{\frac{p}{4-p}}\le c_{\rho,\mu}^N
\leq
\frac{3L_0^2}{2\ell^2}K^{\frac{2p}{p-4}}_{M,N-1}
\mu^\frac{p}{4-p}, \quad \rho\in [1/2,1].
\end{equation}
Therefore, since $\displaystyle\lim_{M \to +\infty}\frac{K_{M,N-2}}{M^{N-1}}= 0$ for $N  =3,\ldots,q+1$, there exists  an integer $M_0 \geq 2$ such that, for every integer $M \geq M_0$ and every $N  =3,\ldots,q+1$, 
\begin{equation}\label{ep-levels1}
\frac{3L_0^2}{2\ell^2}K^{\frac{2p}{p-4}}_{M,N-2}
\mu^\frac{p}{4-p}<\frac{(p-4)2^{\frac{3p-4}{4-p}}}{p}M^{\frac{2p(N-1)}{p-4}}\mu^{\frac{p}{4-p}}.
\end{equation}
For each $N =2,\ldots,q+1$, the monotonicity of $c_{\rho,\mu}^N$ with respect to $\rho\in [\frac{1}{2},1]$ implies that $\{c_{\rho,\mu}^N\}$ is bounded and that $c_{1,\mu}^N\leq c_{\rho,\mu}^N \leq c_{1/2,\mu}^N$. Combining this with \eqref{ep-levels} and \eqref{ep-levels1}, we prove the lemma.
\end{proof}


\begin{proof}[Proof of Theorem \ref{th2}]
Fix $M\geq M_0$ and $0<\mu<\mu_0$, where $M_0 \geq 2$ is defined as in Lemma \ref{prop-windows}. By Lemma \ref{lem-window-lower}, for every $N=2,\ldots,q+1$, we may apply Theorem \ref{minimax} to the functionals $E_\rho(\cdot,\G_M)$. 
From Theorem \ref{minimax}, for every $N=2,\ldots,q+1$, and for almost every \( \rho \in [\frac{1}{2},1] \), there exists a bounded sequence \(\{u_{\rho,\mu,n}^N\}_{n=1}^{\infty} \subset H^1_\mu(\mathcal{G}_{M})\) such that
\begin{equation}\label{5.1}
E_\rho(u_{\rho,\mu,n}^N, \mathcal{G}_M) \to c_{\rho,\mu}^N,
\end{equation}
and
\begin{equation}\label{5.2}
E'_\rho(u_{\rho,\mu,n}^N, \mathcal{G}_{M}) + \lambda_{\rho,\mu,n}^N(u_{\rho,\mu,n}^N, \cdot)_2 \to 0\purple{,}\blue{} \quad \text{in the dual of } H^1(\mathcal{G}_{M}),
\end{equation}
where
\begin{equation}\label{lagrange}
\lambda_{\rho,\mu,n}^N := -\frac{1}{\mu} E'_\rho(u_{\rho,\mu,n}^N, \mathcal{G}_{M})[u_{\rho,\mu,n}^N]. 
\end{equation}
Moreover, there exists a sequence \(\{\zeta_{\rho,n}^N\} \subset \mathbb{R}^+\) with \(\zeta_{\rho,n}^N \to 0^+\) such that, if the inequality
\begin{equation*}
%\int_\mathcal{G} |\varphi'|^2 dx+ \lambda_n\int_\mathcal{G} \varphi^2  dx-(p-1)\rho\sum _{v\in \V}|u_n(v)|^{p-2}\varphi^{2}(v)= 
E''_\rho(u_{\rho,\mu,n}^N, \mathcal{G}_M)[\varphi, \varphi] + \lambda_{\rho,\mu,n}^N  \|\varphi\|_2^2 < -\zeta_{\rho,n}^N \|\varphi\|^2 
\end{equation*}
holds for every \(\varphi\in W_n\setminus\{0\}\), where \(W_n\subset T_{u_{\rho,\mu,n}^N}H^1_\mu(\mathcal{G}_{M})\), one has \(\dim W_n\leq N\).
For each $N=2,\ldots,q+1$, the monotonicity of $c_{\rho,\mu}^N$ with respect to $\rho\in [\frac{1}{2},1]$ implies that $\{c_{\rho,\mu}^N\}$ is bounded and that $c_{1,\mu}^N\leq c_{\rho,\mu}^N \leq c_{\frac{1}{2},\mu}^N$.

By Lemma \ref{lamuna+}, we have $\displaystyle\lambda_{\rho,\mu}^N=\lim_{n \to \infty }\lambda_{\rho,\mu,n}^N\geq 0$. Then,  by Lemmas \ref{lemurhosolve} and \ref{lemurhoconverge}, there exists $u_{\rho,\mu}^{N} \in H^1(\G_M)$ satisfying \eqref{eqpdnomu} and 
\begin{equation}\label{eqpfth2cnrm}    
E_\rho(u_{\rho,\mu}^{N},\G_{M})=\lim_{n \to \infty} E_\rho(u_{\rho,\mu,n}^{N},\G_{M})=c^N_{\rho,\mu}>0.
\end{equation}
\iffalse
We now prove that \(\lambda_{\rho,\mu}^{N}=0\) is not possible either. Argue by contradiction, suppose that \(\lambda_{\rho,\mu}^{N}=0\).
On one hand, if $\K=\emptyset$, where $\K$ is the compact core of $\G_M$ (note that the graphs $\G$ and $\G_M$ share the same compact core $\K$.), then, every edge \blue{$\e \in \E_{M}$} is an unbounded edge, and thus, by \cite[Remark 4 in Section 3.3]{AD}, \purple{there exist $\{a_\e\}_{\e \in \E},\{b_\e\}_{\e \in \E} \subset \R$ such that}\blue{there exist coefficients $a_\e,b_\e\in\mathbb{R}$, $\e\in\E$, such that}
      \begin{equation*}
u_{\rho,\mu,\e}^{N}(x)=a_\e x+b_\e\quad \text{on every edge } \e\in \E,
\end{equation*}
where  $u_{\rho,\mu,\e}^N:=u_{\rho,\mu}^N|_\e$.
Since $u_{\rho,\mu}^N \in H^1(\G_{M})$, we must have $a_e=b_e=0$ for all $\e \in \E$, which implies
$$\displaystyle  \lim_{n \to \infty} E_\rho(u_{\rho,\mu,n}^{N},\G_{M})=E_{\rho}(u_{\rho,\mu}^N, \G_{M})=0,$$
we obtain a contradiction. On the other hand, if $\K\neq\emptyset$, then, by  Lemma \ref{lemamc}, the solution $u_{\rho,\mu}^{N}$ satisfies $E_{\rho}(u_{\rho,\mu}^{N},\G_M) \leq S_{p,\K,\vv_1}$ for all $\rho \in [\frac{1}{2},1]$ and $N=2,\ldots,q+1$.
However, by Lemma \ref{lem-window-lower}, there exists $\mu_0\in (0,\mu_1]$ such that, for every $\mu\in (0,\mu_0)$, every $N=2,\ldots,q+1$, and for almost every \( \rho \in [\frac{1}{2},1] \), 
$$E_{\rho}(u_{\rho,\mu}^{N},\G_M)=c^N_{\rho,\mu} \geq  \frac{(p-4)2^{\frac{3p-4}{4-p}}}{p}M_0^{\frac{2p(N-1)}{p-4}}\mu^{\frac{p}{4-p}} \geq S_{p,\K,\vv_1} +1,$$
which leads to a contradiction.
\fi
We now exclude the possibility that
\(\lambda_{\rho,\mu}^N=0\). Suppose, for contradiction, that \(\lambda_{\rho,\mu}^N=0\). For every edge $\e \in \E_{M}$, by \cite[Remark 4 in Section 3.3]{AD}, there exist coefficients $a_\e,b_\e\in\mathbb{R}$, such that
      \begin{equation*}
u_{\rho,\mu,\e}^{N}(x)=a_\e x+b_\e\quad \text{on every edge } \e\in \E_M,
\end{equation*}
where  $u_{\rho,\mu,\e}^N:=u_{\rho,\mu}^N|_\e$.
Since $u_{\rho,\mu}^N \in H^1(\G_{M})$, we must have $a_\mathcal{H}=b_\mathcal{H}=0$ for every half-line $\mathcal{H} \in \E_M$. In particular, since at least one half-line is attached at each vertex \(\vv_i\in\V_0\), we have
$u_{\rho,\mu}^N(\vv_i)=0$,  $i=1,\ldots,q$.
Testing the equation with \(u_{\rho,\mu}^N\), integrating over every edge and
summing over \(\G_M\), we obtain
\[
\|(u_{\rho,\mu}^N)'\|_2^2
=
\rho\sum_{i=1}^q
|u_{\rho,\mu}^N(\vv_i)|^p
=
0.
\]
Since \(u_{\rho,\mu}^N \in H^1(\G_M)\), it follows that $u_{\rho,\mu}^N\equiv0$ on $\G_M$,
which contradicts \eqref{eqpfth2cnrm}. Therefore, $\lambda_{\rho,\mu}^N>0$. By Lemma \ref{lemurhoconverge}, we conclude that there exists a solution $(u_{\rho,\mu}^{N},\lambda_{\rho,\mu}^{N}) \in H^1(\G_{M}) \times (0, +\infty)$ to the following equation on metric graph $\G_M$ and $u_{\rho,\mu}^{N}$ satisfies $E_{\rho}(u_{\rho,\mu}^{N},\G_{M})=c_{\rho,\mu}^N$ and
\begin{align*}
\begin{cases} 
(u_{\rho,\mu}^N)'' = \lambda_{\rho,\mu}^N u_{\rho,\mu}^N & \text{on every }\e \in \E_M, \\ 
\int_{\G_{M}}\abs{u_{\rho,\mu}^N }^2\, dx = \mu & \\
\displaystyle\sum_{\e \succ \vv} (u_{\rho,\mu,\e}^N)' (\vv)= -\rho|u_{\rho,\mu}^N(\vv)|^{p-2}u_{\rho,\mu}^N(\vv)  & \text{at every } \vv \in \V_0,\\
\displaystyle\sum_{\e \succ \vv} (u_{\rho,\mu,\e}^N)' (\vv)=0  & \text{at every } \vv \in \V \setminus \V_0.
\end{cases}
\end{align*}

For each $N=2,\ldots,q+1$, let $I_N\subset[1/2,1]$ be a full-measure subset such that, for every $\rho\in I_N$, the preceding construction provides a critical point $(u_{\rho,\mu}^N,\lambda_{\rho,\mu}^N)$. Since only finitely many values of $N$ are involved, the set
\[
I_*:=\bigcap_{N=2}^{q+1}I_N
\]
still has full measure. Choose a common sequence $\rho_n\in I_*$ such that $\rho_n\to1^-$. Then, for every $N=2,\ldots,q+1$ and every $n$, $(u_{\rho_n,\mu}^N,\lambda_{\rho_n,\mu}^N)$ solves the $\rho_n$-parameterized version of \eqref{eqpdm} and satisfies
\[
E_{\rho_n}(u_{\rho_n,\mu}^N,\G_M)=c_{\rho_n,\mu}^N,
\qquad
c_{1,\mu}^N\leq c_{\rho_n,\mu}^N\leq c_{1/2,\mu}^N.
\]
By Lemma \ref{lembu}, each sequence $\{u_{\rho_n,\mu}^N\}_n$ is bounded in $H^1(\G_M)$ and is a Palais--Smale sequence for $E(\cdot,\G_M)$ constrained on $H^1_\mu(\G_M)$. Since $N$ ranges over a finite set, a diagonal extraction yields, simultaneously for all $N=2,\ldots,q+1$,
\[
\lambda_{\rho_n,\mu}^N\to\lambda_\mu^N\geq0,
\qquad
u_{\rho_n,\mu}^N\rightharpoonup u_\mu^N\quad\text{in }H^1(\G_M),
\]
where $u_\mu^N$ solves the limiting equation without the mass constraint. If $\lambda_\mu^N=0$, then the argument used above on the half-lines attached to every vertex in $\V_0$ gives $u_\mu^N\equiv0$. Lemma \ref{lemurhoconverge} would then imply $\|(u_{\rho_n,\mu}^N)'\|_2\to0$, and the Gagliardo--Nirenberg inequality
\eqref{eqgn} would consequently yield $E_{\rho_n}(u_{\rho_n,\mu}^N,\G_M)\to0$, contradicting the positive lower bound in \eqref{eqlem55}. Therefore $\lambda_\mu^N>0$. Lemma \ref{lemurhoconverge} now gives
\[
u_{\rho_n,\mu}^N\to u_\mu^N\quad\text{strongly in }H^1(\G_M),
\]
so that $\|u_\mu^N\|_2^2=\mu$ and $(u_\mu^N,\lambda_\mu^N)\in H^1(\G_M)\times(0,+\infty)$ is a solution to \eqref{eqpdm}. Moreover, since $0<\mu<\mu_0$, $M\geq M_0$, and $c_{\rho,\mu}^N$ is monotone in $\rho$, \eqref{sep-levels} yields
$$
E(u^N_\mu,\G_M)=\lim_{n\to+\infty}c_{\rho_n,\mu}^N\in\left[c_{1,\mu}^N,
c_{1/2,\mu}^N\right], \quad N= 2,\ldots,q+1,
$$
and the energies $E(u^N_\mu,\G_M)$, $N=2,\ldots,q+1$, are pairwise distinct. For each \(N=2,\ldots,q+1\), the above construction
therefore yields one pair of normalized solutions $(\pm u_\mu^N,\lambda_\mu^N)$.
Since the index \(N\) takes exactly \(q\) values and the corresponding
solutions have pairwise distinct energies, these solutions form \(q\) distinct
pairs. This gives at least $q$ distinct pairs of normalized solutions.
\end{proof}



\iffalse
\section{Appendix}
Recall that the compact core of $\G$ is defined as the metric subgraph of $\G$ consisting of all the bounded edges of $\G$, and we denote it by $\K$. In the following, we assume that $\G$ has a non-empty compact core $\K$. 

We first recall \purple{some basic definitions of analysis on (unordered) combinatorial graphs}\blue{some basic notions from analysis on (unordered) combinatorial graphs}, see \cite{Gri0} for more details. In what follows, we will also regard $\K$ as \purple{a combinatorial finite graph}\blue{a finite combinatorial graph}. Since $\G$ is connected and $\K$ is non-empty, it follows that $\K=(\V_\K,\E_\K)$ is connected, where $\V_\K=\V$ and $\E_\K$ denotes the set of bounded edges of $\G$. For vertices $v, w\in \V$,   we write $v \sim w$ (i.e. $v$ is \purple{a neighbour of}\blue{a neighbor of} $w$) if $vw \in \E_\K$.

Let $C(\V)$ denote the set of real-valued functions on $\V$. For any function $f\in C(\V)$ and any $\vv \in \V$, the graph Laplacian of $f$ is defined by
$$
\Delta_d f(\vv):=\sum_{\w \in \V, \w\sim \vv}\frac{1}{\ell_{\vv\w}}(f(\w)-f(\vv)),
$$
where each bounded edge $\vv\w \in \E$ is identified with a closed and bounded interval $I_{\vv\w} = [0, \ell_{\vv\w}]$ with $\ell_{\vv\w} > 0$.


For any $p \in[1, \infty]$, we denote by $\|\cdot\|_{\ell^p(\V)}$ the usual $\ell^p(\V)$-norm on $C(\V)$, that is, 
$$
\|f\|_{\ell^p(\V)}:=\left\{\begin{array}{lc}\displaystyle
\Big(\sum_{\vv \in \V}|f(\vv)|^p\Big)^{1/p}, & 1 \leq p<\infty, \\
\displaystyle\sup _{\vv \in \V}|f(\vv)|, & p=\infty .
\end{array}\right.
$$
We define \purple{a semi-norm}\blue{a seminorm} on $C(\V)$ by
$$
[f]_{H^1(\V)}:= \left( \sum_{\vv\w \in \E_\K}\frac{1}{\ell_{\vv\w}}(f(\w)-f(\vv))^2\right)^{1 / 2}.
$$


\purple{We now study the relation between the following two equations: the first one posed on the metric graph $\G$}\blue{We now study the relationship between two problems, one posed on the metric graph $\G$} 
\blue{
\begin{equation}\label{eqpd0}
    \begin{cases} 
u'' = 0 & \text{on every }\e \in \E, \\ 
\displaystyle\sum_{\e \succ \vv} u'_\e(v) = -\rho|u(\vv)|^{p-2}u(\vv) & \text{at every } \vv \in \V_0,\\
\displaystyle\sum_{\e \succ \vv} u'_\e(v) = 0 & \text{at every } \vv \in \V \setminus \V_0,
\end{cases}
\end{equation}
}
\purple{and the second one on the combinatorial graph $\K$}\blue{and the other on the combinatorial graph $\K$}
\blue{
\begin{equation}\label{eqcg}
-\Delta_df=\rho\chi_{\V_0}\abs{f}^{p-2}f \quad \text{in }\V,
\end{equation}
where \purple{$\chi_{\V_0}: \V \to \{0,1\}$ such that $\chi_{\V_0}(\vv):=$}\blue{$\chi_{\V_0}: \V \to \{0,1\}$ is defined by $\chi_{\V_0}(\vv):=$} \begin{cases}
    1 \quad &\text{if }\vv \in \V_0,\\
    0  &\text{if }\vv \in \V \setminus \V_0.\\
\end{cases}$
}
  \begin{lemma}\label{lemamc}
    \blue{\purple{Let $\G=(\V,\E)$ be a noncompact metric graph, $v_\mathcal{H}$ be an end point of a half-line $\mathcal{H} \in \E$, $\rho \in [\frac{1}{2},1]$ and let $u\in H^1(\G)$ be a solution of equation \eqref{eqpd0}.}\blue{Let $\G=(\V,\E)$ be a noncompact metric graph, let $v_\mathcal{H}$ be the endpoint of a half-line $\mathcal{H}\in\E$, let $\rho\in[\frac{1}{2},1]$, and let $u\in H^1(\G)$ be a solution of equation \eqref{eqpd0}.}}  Then $\hat{u}:=u|_\V$ is a solution of \eqref{eqcg} and satisfies
    $$
    \norm{u'}^{2}_2=[\hat{u}]_{H^1(\V)}^2.
    $$
    Moreover, 
    there exists a constant \blue{$S_{p,\K,\vv_{\mathcal{H}}} > 0$},  depending only on \blue{$p$, $\K$ and $\vv_\mathcal{H}$}, such that
    \blue{$$\frac{1}{2}\norm{u'}^{2}_2-\frac{\rho}{p}\sum_{\vv \in \V_0}\abs{u(\vv)}^p=\frac{1}{2}[\hat{u}]_{H^1(\V)}^2-\frac{\rho}{p}\sum_{\vv \in \V_0}\abs{\hat{u}(\vv)}^p \leq S_{p,\G}.$$}
  \end{lemma}
  \begin{proof}
      Let $\rho \in [\frac{1}{2},1]$, and let $u\in H^1(\G)$ be a solution of equation \eqref{eqpd0}. Then, for any $\e \in \E$, we have $u_{\e}''=0$. It follows from \cite[Remark 4 in Section 3.3]{AD} that there exist $\{a_\e\}_{\e \in \E},\{b_\e\}_{\e \in \E} \subset \R$ such that
      \begin{equation*}
u_{\e}(x)=a_\e x+b_\e\quad \text{on every edge } \e\in \E.
\end{equation*}
For any half-line $\e \in \E\backslash \E_\K$, since $u \in H^1(\G)$, \purple{we have that}\blue{we have} $a_e=b_e=0$.
Thus, for any $\vv \in \V$,
$$
\sum_{\e \succ \vv} u'_\e(\vv)=\sum_{\w\in\V, \w\sim \vv}u'_{\vv\w}(\vv)=\sum_{\w\in\V, \w\sim \vv}\frac{1}{\ell_{\vv\w}}(u(\w)-u(\vv))=\Delta_d \hat{u}(\vv),
$$
which implies that $\hat{u}:=u|_\V$ is a solution of \eqref{eqcg}. Moreover, we have
$$
\norm{u'}^{2}_2=\int_\G\abs{u'}^2\,dx=\int_\K\abs{u'}^2\,dx=\sum_{\vv \w \in \E_\K}\frac{1}{\ell_{\vv\w}}(f(\w)-f(\vv))^2=[\hat{u}]_{H^1(\V)}^2,
$$
which implies 
\blue{$$
\frac{1}{2}\norm{u'}^{2}_2-\frac{\rho}{p}\sum_{\vv \in \V_0}\abs{u(\vv)}^p=\frac{1}{2}[\hat{u}]_{H^1(\V)}^2-\frac{\rho}{p}\sum_{\vv \in \V_0}\abs{\hat{u}(\vv)}^p.
$$}
\blue{Let $\mathcal{H}\in \E\backslash \E_\K$ be a half-line of $\G$ and $\vv_\mathcal{H}$ be \purple{an endpoint of $\e_\mathcal{H}$}\blue{the endpoint of $\mathcal{H}$}. It follows from $u_{\e_\mathcal{H}}=0$ and $u \in H^1(\G)$ that $u(\vv_\mathcal{H})=0$.}  Since \blue{$C_{\vv_\mathcal{H}}(\V)$} is a finite-dimensional space, $[\cdot]_{H^1(\V)}$ and $\norm{\cdot}_{L^p(\V)}$ are equivalent norms on  \blue{$C_{\vv_\mathcal{H}}(\V)$}. 
Thus, there exists a constant $C$, depending on \blue{$p$, $\K$ and $\vv_\mathcal{H}$} such that, 
$$
[f]_{H^1(\V)} \leq C\norm{f}_{L^p(\V)}, \quad \forall f \in  \blue{C_{\vv_\mathcal{H}}(\V)},
$$
which implies that
$$
\frac{1}{2}[\hat{u}]_{H^1(\V)}^2-\blue{\frac{\rho}{p}\sum_{\vv \in \V_0}\abs{\hat{u}(\vv)}^p}\leq \frac{C^2}{2}\norm{\hat{u}}_{L^p(\V)}^2 - \frac{\rho}{p}\norm{\hat{u}}_{L^p(\V)}^p \leq \frac{C^2}{2}\norm{\hat{u}}_{L^p(\V)}^2 - \frac{1}{2p}\norm{\hat{u}}_{L^p(\V)}^p\leq \blue{S_{p,\K,\vv_{\mathcal{H}}}}
$$
for some constant \blue{$S_{p,\K,\vv_{\mathcal{H}}}$}.
  \end{proof}
On the other hand, \purple{following the same argument as in Lemma \ref{lemamc}}\blue{by the same argument as in the proof of Lemma \ref{lemamc}}, one can prove that if $f \in C(\V)$ is a solution of \eqref{eqcg}, then the function $\tilde{f} \in H^1(\G)$ defined by
$$\tilde{f}_\e(x):=\begin{cases}
\purple{\frac{u(\w)-u(\vv)}{\ell_{\vv\w}}x+u(\vv)}\blue{\frac{f(\w)-f(\vv)}{\ell_{\vv\w}}x+f(\vv)} &\text{if }\e=\vv\w \in \E \text{ is a bounded edge},\\
    0 &\text{if }\e \in \E \text{ is an unbounded edge},\\
\end{cases}$$
 is a solution of \eqref{eqpd0}, where, for any bounded edge $\e=\vv\w \in \E$, $\vv$ is identified with $0$.
 
Now, for $N \geq 1$, \purple{let $\mathbb{Z}^N:=(\V_N,\E_N)$ be a metric graph}\blue{let the lattice graph $\mathbb{Z}^N:=(\V_N,\E_N)$ be defined}, where 
\[\V_N= \{x=(x_1,...,x_N)\mid x_i \in \mathbb{Z}, 1 \leq i \leq N\},\]
\[
\E_N=\left\{ (x,y) \mid x,y\in\mathbb{Z}^{N},\mathop{\sum\limits _{i=1}^{N}\abs{x_{i}-y_{i}}=1}\right\},
\]
and, for any $\e \in \E_N$, $\ell_\e=1$.  Similarly, if $f \in C(\V_N)$ is a solution of 
\begin{equation*}
-\Delta_df=\abs{f}^{p-2}f \quad \text{in }\V_N,
\end{equation*}
then $\tilde{f}$ is a solution of 
\begin{equation}\label{eqzn}
    \begin{cases} 
u'' = 0 & \text{on every }\e \in \E_N, \\ 
\displaystyle\sum_{\e \succ \vv} u'_\e(v) = -|u(\vv)|^{p-2}u(\vv) & \text{at every } \vv \in \V_N,
\end{cases}
\end{equation}
where $\tilde{f}: \mathbb{Z}^N  \to \R$ is defined by $$\tilde{f}_\e(x):=
\purple{\frac{u(\w)-u(\vv)}{\ell_{\vv\w}}x+u(\vv)}\blue{\frac{f(\w)-f(\vv)}{\ell_{\vv\w}}x+f(\vv)} \quad \text{for any }\e=\vv\w \in \E \text{, where \purple{$v$ is identified with $0$}\blue{the endpoint $\vv$ is identified with $0$ on the edge $\e=\vv\w$}}.
$$
Thus, by \cite[Corollary 4]{hl}, for any $N \geq 3$ and any $p > \frac{2N}{N-2}$, problem \eqref{eqzn} admits a positive solution.
\fi
\subsection*{Conflict of interest}

The authors declare no conflict of interest.

\subsection*{Ethics approval}
 Not applicable.

\subsection*{Data Availability Statement}
Data sharing not applicable to this article as no datasets were generated or analysed during the current study.

\subsection*{Acknowledgements}
 C. Ji is partially supported by the National Natural Science Foundation of China(No. 12571117).
 
 
\begin{thebibliography}{99}
%\bibitem{Ac}N. Ackermann, A nonlinear superposition principle and multibump solutions of periodic Schr\"odinger equations, J. Funct. Anal. 234 (2006), 277-320.

\bibitem{ABD}
R. Adami, F. Boni, S. Dovetta, Competing nonlinearities in NLS equations as source of threshold phenomena on star graphs, J. Funct. Anal. 283 (2022), Paper No. 109483, 34 pp.

\bibitem{ABR}
R. Adami, F. Boni, A. Ruighi, Non-Kirchhoff vertices and nonlinear Schr\"odinger ground states on graphs, Mathematics 8 (2020), 617.

\bibitem{ACFN}
R. Adami, C. Cacciapuoti, D. Finco, D. Noja, Constrained energy minimization and orbital stability for the NLS equation on a star graph, Ann. Inst. H. Poincar\'{e} C Anal. Non Line\'{e}aire 31 (2014), 1289-1310.

\bibitem{AST}
R. Adami, E. Serra, P. Tilli, NLS ground states on graphs, Calc. Var. Partial Differential Equations 54 (2015), 743-761.

\bibitem{AST2}
R. Adami, E. Serra, P. Tilli, Threshold phenomena and existence results for NLS ground states on metric graphs, J. Funct. Anal. 271 (2016), 201-223.

\bibitem{AST3}
R. Adami, E. Serra, P. Tilli, Negative energy ground states for the $L^2$-critical NLSE on metric graphs, Commun. Math. Phys. 352 (2017), 387-406.

%\bibitem{Adams}R.A. Adams, J.J.F. Fournier,  Sobolev spaces, Second edition, Pure Appl. Math. (Amst.), 140, Elsevier/Academic Press, Amsterdam, 2003. xiv+305 pp, ISBN:0-12-044143-8.

\bibitem{AD}
W. Adkins, M. Davidson, Ordinary Differential Equations, Springer, New York, 2012.


\bibitem{AHJ}
C.O. Alves, Z. He, C. Ji, Existence and multiplicity of normalized solutions to a large class of elliptic equations on bounded domains, arXiv:2507.04624, 2025.

\bibitem{BBDT}
D. Barbera, F. Boni, S. Dovetta, L. Tentarelli, Normalized solutions of one-dimensional defocusing NLS equations with nonlinear point interactions, arXiv:2503.21700, 2025.

\bibitem{BBDT2}
D. Barbera, F. Boni, S. Dovetta, L. Tentarelli, Doubly nonlinear Schr\"odinger normalized ground states on 2D grids: existence results and singular limits, arXiv:2511.00550, 2025.

%\bibitem{Bere}H. Berestycki and P.-L. Lions, Nonlinear scalar field equations. II. Existence of infinitely many solutions, Arch. Rational Mech. Anal. 82 (1983), 347-375.
  
\bibitem{Be}
G. Berkolaiko, P. Kuchment, Introduction to quantum graphs, Math. Surveys Monogr., 186, American Mathematical Society, Providence, RI, 2013. xiv+270 pp, ISBN:978-0-8218-9211-4.

%\bibitem{bogachev2007measure}V.I. Bogachev, M.A.S Ruas, Measure theory, volume 1, Springer, Heidelberg, 2007.

\bibitem{BD}
F. Boni, S. Dovetta, Prescribed mass ground states for a doubly nonlinear Schrödinger equation in dimension one, J. Math. Anal. Appl. 496 (2021),  Paper No. 124797, 16 pp.

\bibitem{BD2}
F. Boni, S. Dovetta, Doubly nonlinear Schrödinger ground states on metric graphs, Nonlinearity 35 (2022),  3283-3323.

\bibitem{BDS} F. Boni, S. Dovetta, E. Serra, Normalized ground states for Schr\"odinger equations on metric graphs with nonlinear point defects, J. Funct. Anal. 288 (2025), Paper No. 110760, 40 pp.

\bibitem{Bort}
J. Borthwick, X. Chang, L. Jeanjean, N. Soave, Normalized solutions of $L^2$-supercritical NLS equations on noncompact metric graphs with localized nonlinearities, Nonlinearity 36 (2023), 3776-3795.

\bibitem{Bort0}
J. Borthwick, X. Chang, L. Jeanjean, N. Soave, Bounded Palais-Smale sequences with Morse type information for some constrained functionals, Trans. Amer. Math. Soc. 377 (2024), 4481-4517.

\bibitem{Ha}
H. Brezis,
Functional analysis, Sobolev spaces and partial differential equations,
Universitext, pp. xiv+599, Springer, New York, 2011, ISBN:978-0-387-70913-0.

%\bibitem{Ha2} H. Brezis, P. Mironescu, Where Sobolev interacts with Gagliardo-Nirenberg, J. Funct. Anal. 277 (2019), 2839-2864.

%\bibitem{Bu} B. Buffoni, M.J. Esteban, E. S\'er\'e, Normalized solutions to strongly indefinite semilinear equations, Adv. Nonlinear Stud. 6 (2006), 323-347.

\bibitem{CDS}
C. Cacciapuoti, S. Dovetta, E. Serra, Variational and stability properties of constant solutions to the NLS equation on compact metric graphs, Milan J. Math. 86 (2018), 305-327.

\bibitem{Ca}
P. Carrillo, D. Galant, L. Jeanjean, C. Troestler, Infinitely many normalized solutions of $L^2$-supercritical NLS equations on noncompact metric graphs with localized nonlinearities, Discrete Contin. Dyn. Syst., doi: 10.3934/dcds.2025085, 2025.


\bibitem{Cha}
X. Chang, L. Jeanjean, N. Soave, Normalized solutions of $L^2$-supercritical NLS equations on compact metric graphs, Ann. Inst. H. Poincar\'{e} C Anal. Non Lin\'{e}aire 41 (2024), 933-959.

%\bibitem{Chen} W. Chen, C. Li, Methods on nonlinear elliptic equations, AIMS Ser. Differ. Equ. Dyn. Syst., 4, American Institute of Mathematical Sciences (AIMS), Springfield, MO, 2010, xii+299 pp, ISBN:978-1-60133-006-2, ISBN:1-60133-006-5.

%\bibitem{Cheverry} C. Cheverry, N. Raymond, A guide to spectral theory-applications and exercises, With a foreword by Peter D. Hislop, Birkh\"auser Adv. Texts Basler Lehrb\"ucher[Birkhh\"auser Advanced Texts: Basel Textbooks], Birkhh\"auser/Springer, Cham, 2021. xx+258 pp, ISBN:978-3-030-67462-5, ISBN:978-3-030-67461-8.

\bibitem{Do}
S. Dovetta, Existence of infinitely many stationary solutions of the $L^2$-subcritical and critical NLSE on compact metric graphs, J. Differential Equations 264 (2018), 4806-4821.

\bibitem{Do0}
S. Dovetta, L. Jeanjean, E. Serra, Normalized solutions of $L^2$-supercritical NLS equations on noncompact metric graphs, Proc. Royal Soc. Edinburgh Sec. A: Math., to appear.

\bibitem{Do1}
S. Dovetta, E. Serra, P. Tilli, Uniqueness and non-uniqueness of prescribed mass NLS ground states on metric graphs, Adv. Math. 374 (2020), 107352, 41 pp.

\bibitem{DT}
S. Dovetta, L. Tentarelli, $L^2$-critical NLS on noncompact metric graphs with localized nonlinearity: topological and metric features,  Calc. Var. Partial Differential Equations 58 (2019), Paper No. 108, 26 pp.

\bibitem{DT2}
S. Dovetta, L. Tentarelli, Ground states of the $L^2$-critical NLS equation with localized nonlinearity on a tadpole graph, Discrete and continuous models in the theory of networks, 113-125. Oper. Theory Adv. Appl., 281, Birkh\"{a}user/Springer, Cham.


%\bibitem{Es} M.J. Esteban, E. S\'er\'e, Solutions of the Dirac-Fock equations for atoms and molecules, Comm. Math. Phys. 203 (1999), 499-530.

%\bibitem{evans} L.C. Evans, Partial differential equations, Second edition, Grad. Stud. Math., 19, American Mathematical Society, Providence, RI, 2010. xxii+749 pp, ISBN:978-0-8218-4974-3.

\bibitem{Gn}
S. Gnutzmann, U. Smilansky, S. Derevyanko, Stationary scattering from a nonlinear network, Phys. Rev. A 2011, 83, 033831.

%\bibitem{Gri0} A. Grigor'yan, Introduction to analysis on graphs, University Lecture Series, 71. vol. 8, pp. viii+150, American Mathematical Society, Providence, RI (2018), ISBN: 978-1-4704-4397-9.

%\bibitem{Hof} M. Hofmann, An existence theory for nonlinear equations on metric graphs via energy methods, arxiv No.1909.07856, 2019.

\bibitem{hl}
    B. Hua, R. Li,
    The existence of extremal functions for discrete Sobolev inequalities on lattice graphs,
    J. Differential Equations 305 (2021), 224-241.

%\bibitem{Jeanjean} L. Jeanjean, L. Song, Sign-changing prescribed mass solutions for $L^2$-supercritical NLS on compact metric graphs, arXiv No.2501.14642, 2025.

\bibitem{JPS}
G. Jona-Lasinio, C. Presilla, J. Sj\"ostrand, On Schr\"odinger equations with concentrated nonlinearities, Ann. Phys. 240 (1995), 1-21.

%\bibitem{Kor} E. Korotyaev, Effective masses for zigzag nanotubes in magnetic fields, Lett. Math. Phys. 83 (2008), 83-95.

%\bibitem{KS} V. Kostrykin, R. Schrader, Quantum wires with magnetic fluxes, Comm. Math. Phys., 237 (2003), 161-179.

%\bibitem{Lieb} E.H. Lieb, M. Loss, Analysis, Second edition, Grad. Stud. Math., 14, American Mathematical Society, Providence, RI, 2001. xxii+346 pp, ISBN:0-8218-2783-9.

\bibitem{Ni}
F. Nier, The dynamics of some quantum open system with short-range nonlinearities, Nonlinearity 11 (1998), 1127-1172.

\bibitem{No}
D. Noja, Nonlinear Schr\"{o}dinger equation on graphs: Recent results and open problems, Philos. Trans. R. Soc. Lond. Ser. A Math. Phys. Eng. Sci. 372 (2014), 20130002, 20 pp.

\bibitem{NP}
D. Noja, D. E. Pelinovsky, Standing waves of the quintic NLS equation on the tadpole graph, Calc. Var. Partial Differential Equations 59 (2020), Paper No. 173, 31 pp.

\bibitem{PS}
D. Pierotti, N. Soave, Ground states for the NLS equation with combined nonlinearities on noncompact metric graphs, SIAM J. Math. Anal. 54 (2022), 768-790.

\bibitem{PSV}
D. Pierotti, N. Soave, G. Verzini, Local minimizers in absence of ground states for the critical NLS energy on metric graphs, Proc. Roy. Soc. Edinburgh Sect. A 151 (2021), 705-733.

%\bibitem{Reed} M. Reed, B. Simon, Methods of modern mathematical physics. IV. Analysis of operators, Academic Press [Harcourt Brace Jovanovich, Publishers], New York-London, 1978. xv+396 pp, ISBN:0-12-585004-2.

%\bibitem{Serov} V. Serov, Fourier series, Fourier transform and their applications to mathematical physics, Appl. Math. Sci., 197, Springer, Cham, 2017. xi+534 pp, ISBN:978-3-319-65261-0, ISBN:978-3-319-65262-7.

\bibitem{ST}
E. Serra, L. Tentarelli, On the lack of bound states for certain NLS equations on metric graphs, Nonlinear Analysis 145 (2016), 68-82.

\bibitem{ST1}
E. Serra, L. Tentarelli, Bound states of the NLS equation on metric graphs with localized nonlinearities, J. Differential Equations 260 (2016), 5627-5644.

\bibitem{T}
L. Tentarelli, NLS ground states on metric graphs with localized nonlinearities, J. Math. Anal. Appl. 433 (2016), 291-304.

%\bibitem{W} M. Willem, Minimax Theorems, Birkh\"{a}user, Boston, 1996.
%W. Adkins, M. Davidson, Ordinary Differential Equations, Springer, New York, 2012.
  \end{thebibliography}
  \end{document}
